
%%%%%%%%%%%%%%%%%%%%%%%%%%%%%%%% SCRIPT FOR E-RARA AND MDZ FILES     %%%%%%%%%%%%%%%%%%%%%%%%%%%%%%%%%%%%%%%%%%%%%%%%
%%%%%%%%%%%%%%%%%%%%%%%%%%% fini le 30.04.2025 par F. GOY            %%%%%%%%%%%%%%%%%%%%%%%%%%%%%%%%%%%%%%%%%%%%%%%%
% !TeX TS-program = lualatex
\documentclass{article}
\usepackage[T1]{fontenc}
\usepackage{microtype}
\usepackage[pdfusetitle,hidelinks]{hyperref}

\usepackage{polyglossia}
\setmainlanguage{english}
\setotherlanguages{latin,greek}
\usepackage[series={},nocritical,noend,noeledsec,nofamiliar,noledgroup]{reledmac}
\usepackage{reledpar}

\usepackage{fontspec}
\setmainfont{TeX Gyre Termes}

\usepackage{sectsty}
\usepackage{xcolor}

\usepackage{fancyhdr}
\pagestyle{fancy}
\fancyhf{}
\fancyhead[LE,RO]{\nouppercase{\leftmark}}  
\cfoot{\thepage}
\renewcommand{\headrulewidth}{0.4pt}

% Redefine \section to remove numbering
\usepackage{titlesec}
\titleformat{\section}[block]{\normalfont\scshape\color{gray}}{}{0pt}{} % no number in heading
\titleformat{\subsection}[hang]{\normalfont}{}{0pt}{} % also remove subsection number
\titleformat{\subsubsection}[hang]{\normalfont\footnotesize\color{black}}{}{0pt}{}

% Modify how section marks are stored to exclude numbers
\makeatletter
\renewcommand{\sectionmark}[1]{%
	\markboth{#1}{}} % Only store the section title, without number
\renewcommand{\subsectionmark}[1]{%
	\markright{#1}} % Only store the subsection title, without number
\renewcommand{\numberline}[1]{} % Hide the section number in TOC
\makeatother

\begin{document}

\date{}
        \title{In priorem epistolam Pauli ad Timotheum commemtarii : [Thomas Cajetan], [1531]}
\maketitle
\tableofcontents
\clearpage
\begin{pages} 
\beginnumbering
        
\section*{PRIORIS }
\section*{AD TIMOTHEVM }\pstart missis inuenitur huiusmodi integra salutatio in fine nisi in epistola ad colo. in qua legitar. gratia cum omnibus vobis amen, et non dicitur domini nostri iesu christi etc. ¶ Et apud graecos inuenit subscriptio talis. Ad thessalonicenses secunda scripta athenis. ¶Caietae die. XVII. martii. M.D.XXIX.  \pend
\phantomsection
\addcontentsline{toc}{subsection}{\textit{THOMAE DE VIO CAIETANICARDINALIS SANCTI XISTI. IN PRIOREM EPISTOLAM PAVLI AD TIMOTHEVM COMMENTARII. }}
\subsection*{\textit{THOMAE DE VIO CAIETANICARDINALIS SANCTI XISTI. IN PRIOREM EPISTOLAM PAVLI AD TIMOTHEVM COMMENTARII. }}\pstart \huge\textbf{P}\normalsize AVLVS apostolus iesu christi sum imperium del saluatoris nostri et. Deest. domini iesu christi spei nostrae.] Non solum se apostolum iesu christi explicat sed apostolatum ipsum explicat non quaesitum sed iniunctum sibi a deo saluatore nostro: naturam exprimens diuinam et officium saluandi nos. vt intelligamus iniunctum ei fuisse apostolatum ad exequendum officium salutis nostrae. Nec tantum a deo sed etiam ab homine domino iesu christo iniunctum sibi explicat. Et dicit domini, eo quod ad statum dominii iam peruenerat iesus christus quum iniunxit apostolatum paulo. Adiungit que spei nostrae, eo quod christus imortalis iam et gloriosus sum corpus est spes nostrae futurae bearitudinis. christi nanque resurrectioni innititur spes nostra quod erimus quoque nos immortales sum corpus. TIMOTHEO dilecto. pro. germano.] Eadem dictio haberur hic quae habita est ad philip. dicendo rogo te germane (seu germana )compar. Quadrat autem dictio illa tum aequali tum minori. vnde et ibi dictum est germana compar et hic dicitur germano filio. Significat enim vere ac confone esse talem. vnde germanus filius est qui vere est consonans patri filius. FILIO in fide.P ad differentiam filii naturalis. vt enim patet act. 1o.paulus asciuit sibi timotheum. quem ex his verbis apparet a paulo fuifle conuersum ad fidem:alioquin non appellaret ipsum filium in fide. GRATIA et misericordia et pax.] Superfluunt coniunctiones. Noua salutatio. nam hactenus salutauerat gratia et pax modo tria dicit gratia misericordia pax, ea ratione quod praelatis maxime necessaria est misericordia:tanq̃ loco dei cui proprium est misereri praesint. A DEO patre et Deest. domino iesu christo domino noltro. ] Bis dicit domino. Iemel ablolute domino iesu christoriterum relatiue ad nos domino nostro. Est siquidem absolute dominus omnium, et est peculiariter dominus noster tanquam sponte seruorum.  \pend
\section{[CAPVT I]}\pstart \huge\textbf{S}\normalsize ICVT rogaui te. ] Completa salutatione inchoat tractatum epistolarem. Vbi aduerte quod aduerbium sicut non habet in textu aliquid correspondens ad complendam orationem vsque ad calcem huius capituli, vbi dicitur hoc praeceptum commendo tibi. ita quod illic compler paulus orationem quam hic inchoat. Et est sensus. sicut rogaui te etc. ita hoc praeceptum commendo tibi quemadmodum olim rogaui te nunc tibi commendo hoc praeceptum. Vnde omnia interiecta ab hinc illucusque vere interiecta sunt ad explicandum necessitatem praecaepti et quod olim dederat et quod modo commendat. VT remaneres ephesi quum irem in macedoniam, vt denunciares. pro. vt praeciperes quibusdam ne aliter docerent.] Dictio graeca interpretata aliter docerent composita est ex alio et doctor : varietatem doctorum clare significans. Mandauit paulus timotheo vt praeciperet quibusdam ne essent varii doctores: quod esset variam inducere doctrinam. Et quantum subiuncta insinnant, aduersus illos loquitur qui aliam de legalibus doctrinam inducebant. nam tales doctores explicabit paulo infenus. NEQVE intenderent fabulis.] Thalmudicas fictiones prohibet. ET genealogiis interminatis.] Non dicit neque intenderent genealogiis sed addit interminatis : quae scilicet sunt sine termino in sacris sen autenticis scripturis, multae enim genealogiae tunc temporis erant interminatae. Et rationem huius vltimi subdit. QVÆ quaestiones praestant magis quam aedificationem. pro. dispensationem dei quae est in fide] Interminatae siquidem genealogiae quod praebeant quaestiones clarum est. Et q magis multiplicandis quaestionibus prosint quam diuinae dispensationi seu gubernationi quam exercet deus in fide paulus docet. nam ad dei gubernationem qua gubernat ea quae sunt fidei parum aut nihil prodest notitia huiusmodi genealogiarum : vt patet de multis genealogiis quas hemmus in libris paralipo.  FINIS autem precaepti est di lectio. ] Dixerat vt preciperes quibusdam ne vati essent doctores, perperea vnitatem finis omnium precaeptorum legis  \pend
\marginpar{[ p.140 ]}
\section*{CAPVT }
\section*{.I. }\pstart explicat:vt ex ipsa vnitate finis praecaeptorum varietas arguatur doctorum. Finis inquit praecaepti est dilectio, non qualiscunque sed procedens a triplici fonte. Primus est. DE CORDE puro.] hoc est affectu puro a passionibus et a malicia. ET conscientia bona.] Secundus fons est applicatio scientiae applicatio inquam bona, ad agendum ad cauendum etc. Multi n.hñt applicationem sciae malam: et innumera mala inde sequunt. ET fide non ficta.] Tertius est hic fons vltimo loco scriptus : eo quod duo primi naturae sunt, hic autem tertius reuelarae est gratiae. Et dicit fide non ficta, propter hypocritas in fide: qui fingebant se christianos  et nec erant christiani nec iudaei.  Discurre per praecaepta : et inuenies omnia ordinari ad dilectionem talem. huiusmodi enim dilectio est dilectio dei et proximi ex vera fide in intellectu et ex puro corde in affectu et ex conscientia bona in applicatione et exequutione : quae perficit et mandata primae et secunde tabulae  Haec sunt quae paulus proponit tractanda in hac epistola. Verum antequam sigillatim parriculas propositas tractet, necessitatem commemorat remansionis timothei ephesi ad precipiendum quae dixit ex casu aberrantium ab his. Vnde dicit. A QVIBVS.] praecaeptis et conditionibus finis praecaeptorum. QVIDAM aberrantes. ] Similes pseudapłis significat. CONVERSI sunt in vaniloquium. ] ab vtili doctrina. VOLENTES esse legisdoctores. ] Hinc apparet iudaeos esse istos qui doctrinam legis moysi ingerebant. Et hinc coniicitur istos esse quos appellauit alia docentes. Et iungitur vaniloquio non doctrina legis, sed voluntas vtendi officio doctoris legis versa dicitur in vaniloquium. Et ratio subditur. NON intelligentes neque quae loqurntur.] ipsimet praecaepta legis. NEQVE de quibus affirmant.] obligari ad praecaepta legis.  \pend\pstart ¶ SCIMVS autem quod bona est lex.] Intendit paulus commemorare vt timotheus non solum praecipiat vitari fabulas et quaestiones de genealogiis sed enim doctores legis moysi male applicantes illam: quod esset si applicarent illam ad obligandum iustificatos in christo. Et propter hoc non repraehendit legem sed laudat eam : dicendo scimus autem quod lex bona est. Vituperat autem male applicantes eam subiumgendo. SI QVIS ea legitime vtatur. ] hoc est si ꝗs eam sum ipsammet legem applicet. Et ne vageris quaerendo quod est legitime applicare explicat subiungendo. SCIENTES. pro. sciens hoc.] Ecce legitimus vsus sciens hoc. QVIA. pro. quod lex iusto non est posita ] Formalis est sermo pauli.nam quum lex non superflue ponatur et iustus quatenus iustus non indigeat lege, formaliter venficatur quod iusto non est lex scripta posita. Rursus quum medium legis sit poena et iustus quatenus iustus nulli sit obnoxius poenae, formaliter verificatur quod iusto non est lex pofita.  Et intendit per hoc quod quemadmodum male applicatur lex ad iustos, ita male applicatur ad iustificatos in christo. et propterea subiunget de seipso iustificato in christo. Cum qua veritate stat quod si iustus recedat a iustiriaistarim inuenit in ordine eorum quibus est lex posita: vtpote non iustus. SED iniustis.] E regione iusti vniuersaliter nomniat per partes multa genera iniustor. Et primo iniustos in genere. ET non subditis. pro. et iis qui subdi nolunt, ipiis.] in patriam ac parentes etc. ET peccatoribus. ] a recto recedentibus. SCELERATIS. ] qui prae aliis copia flagitiorum exuberant. ET contaminatis. pro.prophanis.] ea q̃ religionis sunt polluentibus. PATRICIDIS et matricidis.] Clara sunt crimina haec. HOMICIDIS fornicariis  masculorum concubitoribus, plagiariis.] Dictio plagiariis significat fures hoium in mancipia, vel vendenda vel redimenda ab ipsismet captis sen sibi coniunctis. MENDACIBVS et periuris. Superfluir et. et siquid aliud sanae doctrinae aduersatur. ] Ne per omnes species discurrat vnico verbo reliqua compraehendit. QVÆ est. ] Superflunnt hae duae dictiones adiunctae ab aliquo qui restringere voluit sequentem particulam ad sanam doctrinam:quum tamen possit alio referri.nam et pont referri ad id quod dixerat sciens hoc : vt monexplicet sciens hoc. SECVNDVM euangelium gloriae beati dei.] ad significandum vnde scitur haec libertas iusti a lege. s. exeuangelio gloriae beari dei. et ad discernendum vere scientes ab illis qui volunt esse legisdoctores et non intelligunt quod dicant, quam nesciunt euangelium gloriae dei beati. Potest nihilominus reterri ad sanam doctrinam, ad explicandum quae est sana doctrina: quae scilicet est sum euangelium gloriae non hominis sed dei non sum statum praesentem sed beatum. Et apte loquens de lege appellat euangelium euangelium gloriae beati dei, quoniam finis iusti qui vere iustus est est gloria apud deum beatum quam oportet esse gloriam foeliciratis aeternae. QVOD creditum est mihi. ] Ad declarandum exfacto in seipso paulo quod dixerat (videlicet lex iusto non est posita) subiungit qualis modo est in statu euangelii et qualis prius fuerit et vnde etc. Ab officio itaque euangelico inchoat: et vt materiam gratiarum actionis supputat multitudinem beneficiorum sibi collatorum ad enangelicum munus subiungendo.  GRATIAS ago. pro. et gratiam habeo ei qui me confortauit. pro. potentem fecit in christo iesu domino nostro.] Superfluit enim Et gratiam habeo ei domino nostro iesu christo. vel gratiam habeo ei qui fecit me potentem christo iesu domino nostro. non fecit me potentem mundo aut mihi ipsi sed christo: hoc est ad seruitium christi. QVOD FIDELEM me existimauit. Ecce aliud beneficium quod me iudicauit fidelem quod me vsus est vt fideli. PONENS in ministerio.] procul dubio suo ad propagationem regni coelorum. QVI prius blasphemus fui ] Explicatis donis in libertate euangelii, explicat crimina quoque propria in statu legis. Et dicit se fuisse blasphemum propter iniurias quas irrogauerat christo. Agnoscens siquidem christm esse verum deum agnoscit peccatum suum in christm fuisse blasphemiam. ET psequutor.] subaudi ecclesiae christi. ET contumeliosus. pro. et violentus] Non solum persequutus fuerat paulus eccelesiam  \pend
\section*{PRIORIS }
\section*{AD TINOTHEVM }\pstart christi, fed et violentiam intulerat: vt patet in actibus apostolorum. SED. Deficit et misericordiam dei consequutus sum.] Superfluir dei. QVIA ignorans feci in incredulitate.] Romnem ex parte qualitatis peccati reddit quare misericordiam consequutus sit quia ignoranter peccauit in statu incredulitatis. Ita quod ignorans refertur ad vtrunque : tum ad feci tum ad incredulitatem. nam et peccatum incredulitaris ex ignorantia in paulo fuit : et factum blasphemiae persequutionis et violentiae ex ignorantia quoque fuit. In cuius signum quam primum audiuit ego sum iesus nasarenus quem tu persequeris statim dixit domine quid me vis facere: Nonnulla siquidem ratio veniae est ex ignorantia peccare: quoniam affectus peccantis ex ignorantia  non tendit per se ad malum: vtpote qui non tenderet in illud si sciret esse malum. SVPERABVNDAVIT autem gratia domini nostri.] Non solum misericordiam consequutus sum condonantem mihi peccata mea, sed etiam superabundauit gratia domini nostri: vtpore promouens me ad apostolatum. CVM fide et dilectione quae est in chro iesu. ] Intellige superabundauit referri non solum ad gratiam, sed etiam ad dona fidei et dilectionis in christo. haec enim tria abundanter in paulo apparent. nam et gratia apostolica in eo exuberat: et fides valde inconcussa tum in perpessis tum in miraculis apparet:dilectio quoq ranta vt ipse dicat. quis infirmatur et ego non infirmor quis scandalizatur et ego non vror etc.   \pend\pstart FIDELIS sermo et omni accaeptione dignus.] Ne libertatem iusti ex seipso tantam ex quodam singulari astruere videatur, subiungit vniuersalem causam huiusmodi libertatis aduentum christi. Et sermonem huius causae communis statui euangelico praemitrit esse tum fidelem hoc est veracem, tum omni accaeptatione dignum: vt tam parti cognoscitinae congruere significetur ex veracitate quam parti affectiuae ex accaeptabilitate. QVIA. pro. quod christus iesus venit in hunc mundum peccatores faluos facere.] Superfluit hunc. Ecce sermo fidelis et omni accaeptione dignus. Et dicendo venit in mundum diuinam iesu chrsti personam significat. Quum enim esset supra mundum, ingressus est mundum factus est enim mundi pars quum factus est homo.¶ Et dicendo peccatores saluos facere explicat quorsum venerit. E regione ad legem (quae posita est ad puniendos peccatores) dicitur quod iesus christus venit ad saluandum peccatores. quod est eripere eos a subiectione legis. QVORVM primus ego sum. ] Relatiuum quorum non refert peccatores absolute, sed peccatores saluatos a christo. Et est sensus. quorum peccatorum qui saluantur per christi aduentum ego sum primus : hoc est precipuus. pctonrum siquidem tunc temporis saluatorum per christum duo fuerunt genera. Alterum peccantium ex infirmitate: vt mulier q̃ erat in ciuitate peccatrix vt petrus qui ex timore negauit et similes. Alterum eorum qui per ignorantiam crucifixerunt christm qui postea adueniente spiritu sancto poenitentiam egerunt : vt dicitur in actibus apostolorum. Paulus his omnibus saluatis peccatoribus prestat, tum quam et blasphemus et psequutor et violentus fuerat. tum quia nec per aduentum spiritus scti et predicationem apostolorum conuersus est vt alii. Et propterea merito inter saluatos peccatores primum (hoc est precipuum) se dicit. ET ideo misericordiam consequutus sum: vt in me primo ostenderet iesus christs oemm patientiam pro. longanimitatem, ad informationem eorum qui credituri sunt illi. pro. in illum in vitam aeternam. ] Primo non est aduerbium sed nomem. et significat precipuum. Reddiderat prius romnem consequutae prius misericordiae ex parte qualitatis peccati sui : modo aliam romnem reddit ex parte exempli tribuendi aliis peccatoribus. Ideo inquit misericordiam consecutus sum vt in me primo (hoc est precipuo) monstraret iesus christs oemm longanimitatem propriam quam seruat expectando peccatorum pnĩam (animi enim in longum propensi est hmoni expectatio) ad statuendum me exemplar aliis qui credituri sunt. Informationem intellige exemplari forma. Et est sensus. vt in me praecipuo ostenderet longanimitatem ad exemplarem informationem itineris in vitam aeternam respectu eorum qui credituri sunt in christm. In ipreo siquod em paulo habemus prae oculis exemplar eundi in vitam aeternam quamuis peccatores simus. Modus autem eundi explicatur per fidem in christum. Et haec est propria ratio propter quam paulus precipuus inter peccatores saluatos misericordiam consecutus est. ¶ REGI autem saeculorum. ] Superius dixerat et gratiam habeo : mon exercet gratiarum actiones, dicendo. regi autem non vnius saeculi sed saeculorum. IMMORTALI. ] ad differentiam regum longaeuae vitae : puta noe qui et saeculo añ diluuium prefuit et saeculo post diluuium, qui tñ non fuit immortalis. INVISIBILI.] Hoc ad differentiam coelestium corporum dicitur coelum. n. licet regat saecula et sit immortale, est tum visibile. SOLI deo.] Deficit sapienti. legendum est. SOLI sapienti. ] suapte natura. reliquae enim substantiae immortales ac inuisibiles quum in nudis naturalibus desipere possint peccando, non sunt suapte natura sapientes: sed quaecunque sunt sapientes ex liberali dono dei sapiunt. Quum dixisset tot regis conditiones, par fuit vt plenae sapiae meminerit : explicando non solum quod est sapiens sed quod est solus sapiens. DEO. ] Faciendum est punctum post soli sapienti, et deinde subiungendum deo. quadrat siquidem post tot adiectiua vt natura ipsius regis explicetur, dicendo deo. HONOR et gloria in saecula saeculorum amen. ] Haec de ĩteriectis. Vnde et rediens ad principium epistolae sicut rogaui te vt remaneres ephesi vt praeciperes quibusdam etc. subiungit. HOC precaeptum commendo tibi fili timothee.] Illo mon quo rogani te vt remaneres et preciperes ephesi : nunc quoque commendo tibi hoc precaeptum de fine omnium praecaeptorum explicato de libertate iustorum a lege de dilectione ex corde puro et fide non ficta et conscientia bona. SECVNDVM praecaedentes in te propherias vt milites in illis bonam militiam.] In actibus apostolorum scribitur quod de timotheo multa bona  \pend\pstart teddebantur testimonia quum paulus assumpsit eum secum. plus autem explicatur hic dicendo que   \pend
\marginpar{[ p.141 ]}
\section*{CAPVT }
\section*{.II. }\pstart praeceflerunt prophetiae de timotheo. Et quantum ex his coniicitur, ab aliqua seu aliquibus spiritualibus personis praedicta fuerant multa bona spiritualia timotheo accessura, quae tam paulo quam timotheo notaerant. et propterea animum timothei accendit ex talibus praecedentibus prophetiis. Er est ordo literae. vt fm praecedentes in te (seu super te seu ad te) prophetias milites in illis (prophetiis: hoc est in prophetaris per illas prophetias) bonam militiam agas bonum militem. HABENS fidem et bonam conscientiam.] Applicat ad ipsum timotheum quae in genere dixerat de fide et conscientia bona, vt timotheus sit primus ibi habens fidem et conscientiam bonam: vtpote qui praesidebat ephesi. QVAM quidam repellentes. pro. qua quidam repulsa.] Et relatiusi qua refert bonam conscientiam: ad ostendendum quanti pendenda est bona conscientia hoc est bona applicatio scientiae. CIRCA fidem naufragauerunt, pro. naufragium fecetunt.] Ecce periculum ac damnum non se uatae bonae conscientiae inde siquidem iacturam incurrerunt fidei. Male siquidem applicando scientiam incurritur error in fide : et sic paulatim perditur fides. EX QVIBVS est hymenaeus et alexander] Nihil inuenio autentice scriptum de his duobus qui repulsa bona conscientia perdiderunt fidem. QVOS tradidi satanae.] corporaliter vexandos: vt de corinthio incestuoso explicatur. VT discant non blasphemare.] Quemadmodum corinthium exconicando tradidit satanae sum carnem vt spiritus saluus fieret, ita istos exconicando tradidit satanae vt discant huiusmodi poenis non blasphemate. Insinuat enim hinc quod iesum christum blasphemabant isti perdita fide. Et forte isti erant doctores docentes aduersus christum postquam a fide ceciderunt. et propterea dicit vt discant non blasphemare.  \pend
\section{CAPVT .II.}\pstart \huge\textbf{O}\normalsize BSECRO. pro. adhortor igitur primum omnium.] Tractata necessitate quare reliquerat timotheum ephesi incipit prosequi proposita. Vnde tanquam ad proposita exequenda rediens illatiue dicit adhortor igitur. Proposuerat autem quod finis praecaepti est dilectio de corde puro etc. et propterea ab executione dilectionis erga omnes inchoat. Primus antem et communissimus dilectionis actus erga omnes est oratio. et ideo tractans illam parriculam dilectio hortatur ante omnia orari pro omnibus. FIERI obsecrationes. pro. preces.] petitorias. ORATIONES. ] hoc est mentis eleuationes in deum. POSTVLATIONES. pro. intercessiones. ] Quae differunt a precibus in hoc quod preces respiciunt rem petitam: intercessiones autem mediae sunt ad reconciliandum homines deo. GRATIARVM ACTIONES.] pro beneficiis accaeptis. PRO omnibus hominibus.] Ecce extensio dilectionis ad omnes homines siue christianos siue non christianos. nam non dicit pro omnibus fratribus sed omnibus hominibus. PRO regibus] qui tunc non erant christiani. Vult pro omnibus hominibus in genere pro regibus autem in specie orari. Nec solum pro regibus sed ET OMNIBVS qui in sublimitate constituti sunt. ] Omnes temporales dominos compraehendit generali nomine fublimitaris seu excellentiae. VT QVIETAM et tranquillam vitam agamus] Non concordant omnino interpretes circa haec vocabula. Verum parui refert tranquillam et quietam vitam an placidam et quietam seu tacitam dixeris. Id tamen aduertendum est quod quum dixisset tum pro omnibns homnibus tum pro regibus et dominis, prius more suo declarat postremum. s. quare pro regibus et in celsitudine constituris vt quietam et tranquillam vitam etc. Siquidem sermo pauli formalis est quum dicit pro regibus et temporalibus dominis orandum: videlicet quatenus regunt et dominantur vt regimen eorum iustum sit. nam inde sequitur iste fructus vt tranquillam et quietam vitam agamus christiani. IN OMNI pietate et castitate. pro. honestate.] Ne quietem et tranquillitatem vitae apperant christiani absolute sed relatiue ad virtutum opera adiungit virtutem pietatis erga parentes patriam etc. et vniuersae morum grauiratis. Et dicit in omni, vt nulla pars pietatis nulla pars honesti moris praetermittatur.  \pend\pstart HOC ENIM bonum est.] Pronomen hoc demonstrat orare pro omnibus hominibus. Et quod primo dixerat (.s. pro omnibus hominibus orandum) manifestat hoc esse bonum in se. ET accaeptum coram saluatore nostro deo. ] Non solum est in se bonum sed gratum deo. Quem nominat saluatorem, quia ex officio saluandi manifestat intentum subiungens. QVI omnes homines vult saluos fieri et ad agnitionem veritaris venire.] Non est hic speculatiue theologe sermo de voluntate beneplaciti : de qua scriptum est, omnia quae voluit dominus fecit in coelo et in terra et quod voluntati eius nullus porest resistere, vt clare patet ex eo quod videmus non omnes homines ad agnitionem veritatis venire. Sed est sermo de voluntate signi : qua deus proponit omnibus hominibus praecaepta salutis doctrinamque euangelicant et illis imputatur qui nolunt venire ad veritaris cognirionem et nolunt salui fieri. de qua scriptum est, quotiens volui congregare filios tuos et noluistis. Ex hoc enim quod deus salutem et agnitionem omnibus vult proponendo, monstratur gratum illi esse et pro omnibus oremus, et simul insinuatur quid est petendum orando pro omnibus videlicet vt saluentur. ¶ VNVS enim deus. ] Manifestat quod deus vult omnes homines saluos fieri et ad agnitionem veritatis venite, tum ex vnitate dei.ex hoc enim quod non nisi vnus est deus omnium hominum curam illi vm incumbere manifestatur: et quum sit natura bonus, consequens est vt omnibus proponat salutem et agnitionem veritatis. Si namque essent plures dii, cogitari posset quod vnus deus haberet curam saluandi aliquos homines et alius haberet curam saluandi alios: sed vbi vnus tantum est deus, illi enim incumbit cura omnium. Tum ex vnitate mediatoris subiugendo. VNVS mediator dei et hominum homo christus iesus.] Si plures essent mediatores inter deum et hominem, existimaretur quod vnus esset mediator pro quibusdam de alius pro aliis sed  \pend
\section*{PRIORIS }
\section*{AD TIMOTHEVM }\pstart ex quo vnus est mediator dei et hominum ad reconciliandum homines deo, illi enim incumbit mediare inter deum et omnes homines. Et propterea idem sequitur.s.quod deus vult omnes homines saluos fieri et ad agnitionem veritatis venite ex hocipso quod vnum constituit mediatorem dei et hominum. Qui tanto magis omnium hominum mediator est ad deum quanto ipse est homo quanto ab ipso nihil humani alienum est. QVI dedit redemptionem semetipsum pro omnibus. ] Ab effectu manifestat quod etiam mediator vult omnes homines saluos fieri et ad agnitionem veritatis veniretex hocipso quod dedit semetipsum precium ad redimendum non aliquos sed omnes. Quod verissimum est non solum quantum ad sufficientiam sed etiam quantum ad efficientiam: pro quanto omnes redemit a morte corporis resuscitando omnes in nouissimo die. CVIVS testimonium temporbus suis confirmatum est.] Superfluunt tres dictiones; cuius confirmatum est. Legendum siquidem est apposiriue. dedit redemptionem semetipsum pro omnibns appositiue testimonium temporibus suis. Vere testimonium diuinae voluntatis quod vult omnes homines saluos fieri et ad agnitionem veritatis venire dedit christus dando seipsum redemptionem pro omnibus. dedit autem hoc testimonium temporibus propnis definitis a diuina sapientia. Et addidit hoc, ne quisquam diceret tardum fuisse testimonium respectiue ad eos qui praecesserant aduentum christi. omnis enim tam tarditas quam anticipatio excludirur dicendo tperibus suis. IN QVO. pro. in quod.] videlicet testimonium promulgandum mundo. POSITVS sum ego predicator pro. praeco. ] Officium preconis est publicare: et paulus publicabat testimonium datum a christo. ET apostolus.] Ne vsurpatum aut vile praeconis officium putaretur, adiunxit et apostolus proculdubio christi. ¶ VERITATEM dico. ] Deficit. IN christomnon mentior.] Parenthesis est haec, explicans officium praeconis et apostoli paulum agere vtranque complectendo partem : videlicet non mentiri et veritatem dicere christo teste. vt nulla falsitas mixta verirati intelligatur. DOCTOR gentium in fide et verirate.] Non solum praeco et apostolus sum veritatis, sed sum specialiter doctor gentium in fide  non in pħia, sed non in fide vana sed et veritate. Doceo gentes veritatem non figmenta : verum non in euidentia scientifica sed in fide. et haec omnia ago ad publicandum testimonium quod dedit christus.  \pend\pstart \huge\textbf{V}\normalsize OLO ERGO viros. ] Tractat illam particulam de corde puro. Continuat autem exercens officium doctoris gentium. ORARE in omni loco, leuantes puras manus sine ira et disceptatione.] Dixerat dominus si offers munus tuum ante altare et recordatus fueris quod frater tuus habet aliquod aduersum te etc.  extendit paulus ad omnem locum idem praeceptum. Docet siquidem quod non solum ante altare sed in omni loco in quo ꝗs orauerit oportet leuari manus mundas sine ira et disceptatione. Nec est sermo de membris corporeis sed de operibus. impedimenta siquidem orationis vbicumque oretur, impuriras operdĩ ira et contentio sunt  Et excluduntur ad literam duo crimina. Alterum auariciae occupantis, dicendo puras manus. communi siquidem vsu occupantes aliena dicimus non habere manus puras. Alterum discordiae cum proximo: dicendo sine ira intus, et disceptatione extra. haec enim duo purganda sunt ante orarionem, auaricia restituendo et discordia reconciliatione. Et meminit mundiciae ab his criminibus tantum: quia reliqua crimina interna poenitentia purgantur, haec autem exigunt etiam externam restitutionem et reconcilianonem. ¶ SIMILITER et mulieres in habitu. pro amictu ornato cum verecundia et sobrietate. pro. modestia, ornantes, pro ornate seipsas.] Quemadmodum a viris abstulit impedimenta orationis tanquam propria viris, ita a mulieribus tollit impedimenta orationis quae sunt velut propria mulieribus. Vnde quod dicit in amictu ornato, sic dicit quemadmodum dixerat in omni loco. Vtrumque enim recensetur ad explicandum modum seruandum a viris quidem in quocumque loco orauerint: a mulieribus vero in quocumque amictu omnato. Et explicatur modus dicendo cum verecundia et modestia ornare seipsas. Vtraque pars valde congruit muueribus: ne ornatus sit inuerecundus, et similiter ne sit imoderatus. Vnde immoderationem ornatus explicando subiungit. NON in tortis crinibus aut auro aut margartis.] Primiriuae ecclesiae tempori congruebat hoc praeceptum:eo quod tunc mulieres quae se christianas profitebantur omnia huius modi a se abdicabant tanquam superflua et vana. Vnde hodie quum multae principes et reginae quibus sum statum congruit ornare se auro et margaritis christianae sint, non contra pauli praeceptum his vtuntur. VEL veste preciosa pro. sumptuosa.] Dixerat in amictu omnato et declarauerat imoderatum omatum:declarat modo amictum, inhibendo vestem sumptuosam. Et quod temporale fuerit huiusmodi praeceptum aurea vestis ceciliae testatur. SED QVOD. pro. quae decet mulieres promittentes. pro. profirentes pietatem per opera bona.] Relatiuum quae ad vestem refertur. et vult vestem cem talem quae deceat mulieres profirentes per bona opera (hoc est profitentes non tantum vetbis sed factis) pietatem:hoc est cultum diuinum. Dictio siquidem graeca non significat pietatem in genere, sed eam quae est erga deum. Perpende ex his omnibus impedimenta ora tionis in mulienbus vbicumque orent esse contraria praedictorum: vtpote appetitum muliebrem deprimentia ad haec vana ne eleuetur illarum mens in deum. ¶ MVLIER in silentio discat cum omni subiectione.] Postquam cor purum tam in viris quam in mulieribus declarauit relatiue ad orationem, declarat in specie puritarem cordis in muliere relatine tum ad disciplinam et doctrinam tum ad dominium. Et qm̃ mulier eget vt docearur, modum explicat tum quo ad linguam dicendo in silentio. Non vult quod in ecelesia interroget doctorem docentem, sed quod audiat quidem et discat silendo:si quid autem restat dubii virum domi interroget. Tum quo ad reuetentiam cum omni subiectione interna et externa.  DOCERE autem ] supple publice. MVLIERI  \pend
\marginpar{[ p.142 ]}
\section*{CAPVT }
\section*{.III. }\pstart non permitto.] Et diximus publice ad differentiam doctrinae domesticae : qua mater debet erudite filios et filias. NEQVE dominari pro. neque autoritatem habere in virum. ] Directe hoc respicit vxores : quas nefas est autoritarem habere in maritos. SED essein silentio.] vt non solum opere sed nec sermone autoritatem sibi vsurpet in virum. ADAM enim primus formatus est, deinde eua.] Exordine quo moyses narrat formationem adae et euae, probat virum praeefse mulieri et non econuerso. ET ADAM non est seductus, mulier aurem seducta in praeuaricatione fuit. ] Eadem vtitur historia : in qua mulier ipsa professa est se seductam a serpente. adam vero non se excusauit vt decaeptum sed vt allectum ab vxore : vt ibi patet. Praestat itaque vir mulieri tum ordine formationis tum vigore rationis. SALVABITVR autem per filiorum generationem si permanserit, pro. si manserint in fide et dilectione et sanctificatione cum sobrietate pro. modestia.] Ordo literae est. saluabitur autem si manserint (ipsa et vir) per filiorum generationem in fide etc. ¶ Meritorium salutis astruit esse coningium mulieri : ex hoc quod si ob filiorum generationem manserit simul cum viro in fide etc. saluabitur ipsa mulier. Et per hoc intellige explicari proprium munus sexus. Non docet paulus salutem mulieris pendere a generatione filiorum (quum facilius saluetur virgo quam coningata) sed docet quod vtendo officio sexus habet vnde ex ipso vsu sexus saluetur : hoc est quod ad salutem eis operatur si ob filiorum generationem manserint cum viris suis in fide etc. Et si vis oemm scrupulum ab hac lectione excludere, ordina saluabitur pro predicato totius enunciationis : ordinando literam sic.si permanserint autem per filiorum generationem in fide et dilectione et sanctificatione cum modestia saluabitur. Sic enim clare significatur quod non pendet salus mulieris a filiorum generatione, sed quod acquiritur ei salus si sic manserint.¶ Et nihilominus aduerte quod non de muliere sed de vxore est sermo: de qua dixerat neque autoritatem habere in virum. de ipsa enim subiungit saluabitur asit etc.  Et meminit specialiter sanctitaris ac modestiae : quoniam tam mundicia in vsu coningii quam moderatio in vsu rerum et gubernatione familiae peculiaris esse debet coniugatis.  \pend
\endnumbering\beginnumbering\section{CAPVT .III.}\pstart \huge\textbf{F}\normalsize IDELIS sermo. ] Prosequitur paulus particulam de corde puro quo ad constitutos in ecclesiasticis ministeriis. Praeponit aurem huic tractarui fidelis sermo ad significandam veritatem dicendorum : ne quis imputaret paulo quod nimis exigeret ex proprio capite in episcopis et diaconis. Totum hoc excluditur dicendo fidelis sermo. significatur enim quod est fermo fidelis domino, quod est iuxta fidem quam debet paulus domino et non sum propriam sententiam. SI QVIS episcopatum desiderat bonum opus desiderat.] Initium fidelis sermonis est si quis superintendentiam cupit, bonum (seu praeclarum) opus desiderat. Non dicit si quis dignitatem si quis gradum si quis prouentus si quis honores si quis gloriam episcopatus : sed si quis ipsum episcopatum (hoc est ipsam superintendentiam ipsum superintendere) desiderat, proculdubio desiderat opus quum superintendere sit opus, non qualecunque sed opus bonum ex suo genere immo opus praeclarum ex proprio genere. Nec hoc eget probatione aliqua : quoniam manifeste constat sic esse. Nullum itaque peccatum est ex proprio genere desiderare episcopatum: hoc est ipsum superintendere. Secus autem est de appetitu gradus honoris prouentuum et huiusmodi. nihil enim horum est opus bonum: quum nec sit opus. OPORTET enim. pro. igitur  episcopum irrepraehensibilem esse.] Ex ratione excellentis bonitatis huiusmodi operis quod est superinrendere illatiue dicit oporter igitur episcopum esse irrepraehensibilem. Et haec est prima conditio episcopi relatiue ad malos mores vt nihil in eo sit repraehensibile. Et vere ex eo quod aliis superintendit repraehendendis, manifeste infertur quod oportet ipsum esse irrepraehensibilem. qua enim fronte porest alios repraehendere qui repraehensibilis est. diceretur namque ei medice cura teipsum. VNIVS vxoris virum.] Secunda conditio negatiue intelligenda est : hoc est non plurium vxorum virum. Non enim exigit paulus ab episcopo quod fuerit aut sit coniugatus (qui consulit meliorem esse coelibatum scribendo ad corin.) sed exigit quod si habet vrorem sit vir vnius vxoris. Sic enim litera sonat dicendo esse. Quo contra de vidua inferius dicet in ca. 5. non praesenti sed preterito vtens tempore quae fuerit vnius viri vxor. Et dixit hoc, quia tunc temporis multi plures vxores habebant imitantes patres veteris testamenti : quum nullibi legamus hoc prohibitum. SOBRIVM pro. vigilem.P Tertia conditio vigilantiam summe necessariam supintendenti explicat. PRVDENTEM pro. modestum.] Quarta conditio summe quoque est necessaria superintendenti. vt vigilantia comitem habeat moderationem : vt ipse sic sit vigil vt tñ sit modestus et sic modestus vt tñ sit vigil. ORNATVM.] proculdubio compositione gestuum actionum eloquii et aliis bonis moribus. PVDICVM.] Superfluit pudicum. Sexto enim loco dicitur. HOSPITALEM.] Rara est haec hodie virtus. DOCTOREM.] non laurea, sed scientia sufficiente ad docendum. NON vinolentum. pro. non vinosum.] hoc est deditum vino. NON percussorem.] Quamuis paulus de iniusta pcussione loquatur, turpe tñ est epreo eriam iuste percutere, debet .n. aliena non propria manu percutere delinquentes.  Post non percussorem deficit enim particula. s. NON tutpilucrem. ] Haec est enim decima conditio in textu pauli : clare sonans quod episcopus alienus esse a turpi lucro debet. Quod si seruaretur, non audirentur lucra tam turpia vt turpe sit ea dicere. Nec loquor de simonia, sed proprie de turpi lucro sparso nimis in episcopis ecclesiarum christi. SED modestum. pro. sed mitem.] Esse mitem plus est quam esse mansuetum mitis enim non solum iram temperat, sed dulcedinem quandam adhibet. NON litigo\pend
\section*{PRIORIS }
\section*{AD TIMOTHEVM }\pstart sum.) Valde repugnat officio supintendentis aliisivt ipse litibus incumbat. alienus. Hesse debet ab ol pugna tam verbog q̃ factor. Sola sigdem pugna fidei et pietatis ei sufficit. NON cupidum. pro.non auarum. ] Non est sermo de quacuque cupiditatejsed de cupiditate argentiiauri et hmoni: quam auaritiam di cimus.V SVE domui bum praepositum, pro. praesidentem P hoc est ppriam domum bene regentem, vult.n. paulus experientia comprobatsi fuisse q is qui fit episcopus sit idoneus ad bene regendum ecclesiam:expenentiam autem sumi ex hoc q bene regit propriam domsi. FILIOS habentem subditos cum oi castitate. pro. reuerentia.p seu honestate. Ne intelligas quod oporteat episcopum habere filios (qm̃ hoc paulus non dicit) sed quod episcopus habens filios habeat illos subditos cum oi reuerentia seu honestate. Ex hac.n.experientia tenendi filios subditos cum oi honestate seu reuerentiaispes gignit quod tenebit christianos ac clericos subditos cum oi honestate seu io uerentia sicut rexit filios. SI QVIS autem domui suae pesse nesciit, qumo ecclesiae dei diligentiam habebit pro. curam aget?b Ron clara, solido innixa fundamento. NON neophytum.P Graecam dictionem religt inrerpres: qusi potuisset dicere non nouitium in ministerio ecclesiastico vel in fide. Et romnem reddit. NE IN superbiam elatus in iudicium incidat diaboli. P Occasio supbiendi est si gs noniter comnersus ad fidemsseu non initiatus in ministeriis ecclesiae, praeferat ad supintendendum ecclesiae: et sic merito superbiae incidat in damnationem similem damnationi diaboli, qui ppter supbiam damnatus est. Loqui auĩt paulum de nouirio in ministerio ecclesiaejex subiunctis de pmotione diaconog satis patet.nam si illos vult primum proban, multo magis promouendos in epreos. POPORTET autem illum et testimonisi bonum habere ab iis qui fons sunt. P Vltima condirio episcopi est vt sit bonae famae eriam apud infideles. Er romnem fubdir. VT non in oppro brium incidat. D hoc est vt non incidat in vituperationem, vituperaret. n. si effet apud eos malae famae. ET in laqueum diaboli.P Dictionem graecam religt interpres:forte rectius dixisset calumniatoris. Nam si malae famae est apud infideles, laqueis exponitur hominum calumniatogex ipsa mala fama audentium ctiam quae non feeit imputare, et illaqueare eum detractionibus jaccusationibus et alus huiusmodi.  \pend\pstart CDIACONES. pro. diaconos similiter pudicos. pro. compositos.P Dictionem quog graecam interpres reli quit diacones: quum potuisset latine dicere ministros. Et qum non legimus alios in sacra scriptura diacor nos q̃ illos in actis aplog institutos ministros tpalium ecclesiae, ideo de eisdem intelligo paulum loqui. nam etiam ephesi fuisse diaconos dispensationis tpaliumipatet ex subinnctis: vbi dicet ne grauct ecclesia viduis quibusdam. Vbi et aduerte non loqui paulum de ordinibus sacris, sed de officiis. vñ nullam mentionem facir sa cerdotis, sed epi et diaconi: quorum vtrunque nomen est officii. illud supintendendilhoc ministrandi. P Pri mam itaque conditionem ministrorum ecclesiae, explicat compositionem in moribus in actibusun gestibus. NON bilingues. P duplici vtentes sermone:mon vnum mon contrarium enim sine pluribus dicendo. NON multo vino deditos. P Non phibet eis vsum vini sed multsi. NON turpe lucrum sectantes. P. Non solum ab illicito sed enim aturpi lucro vult esse alienos. sunt. n.multa lucra licita quae tñ sunt turpia. HABENTES mysterium fidei in conscientia pura.per Non sufficit vt diaconi hemant fidemjsed exigit ab eis vt habeant intimum fidei arcanum. Quod consistit in conformitate animi ad fidemjin conformitate volunt aris et affectus ad fidem. hoc est. n.intimum fidei. Et hoc non est habendum qualitercumque sed in conscientia purajsed in appli catione scientiae pura ab oi errore et passione oportet. n.ministros ecclesiae et intus habere animum confor. mem fideii et ministrando applicare pure scienriam ad opus.VET hi autem, pro. quidemiprobent primum: et sic] Loco harum duarum dictionum legendum est. DEINDEj ministrent. Deficit. sic, nullum crimen habentes.P De epreo dixerat non nouitium: de diaconis autem dicit quod primum examinent jexperientia oper discernant. postea ministrent sic. s. probati quod nullum habeant crimensquod nullus possit illos criminari.P MVLIERES. pro.vro res.E proculdubio eporum et diaconorum. de aliis siqdem mulieribus superius locutus fuerat.SIMILITER pudicas, non detrahentes. pro. non calsiniatrices, sobrias, pro. vigiles, fideles in oibus.P Conditiones cla rae sunt.V DIACONES. pro. diaconissint vnius vxoris viri.P Intellige hoc sicut de epreo expositum est. QVi filiis suis bene praesint.P Supfluit suis. ET suis domibus.P Exigit a pmouendis in diaconos similem ex perientiam circa regimen domus et filiog qualem exegerat a pmouendo in episcopum.V QVI.n. bene ministraue rint.2 Post conditiones regsitas ad diaconos fructum eorum subiungit. Nec dicit q.n. ministrauerintised ꝗ bene ministrauerint. GRADVM bonum sibi acqrent. pro. acqrunt.P hoc est altiori officio dignos se consti tuunt. hic.n. fructus suapte natura sequit bonam ministrationem inferioris officil. ET multam fiduciam in fide quae est in christo iesu. P Actiue et passiue verificat quod acqrunt sibi multam fiduciam in fide iesu christi. Actine quidem: quia bum ministrando, magis ac magis confidunt et audent in iis quae sunt fidei: vt patet de stephano. Passiue vero: quia bene ministrando acgrunt vt alii multum confidant in ipsis in iis quae sunt fidei.  \pend\pstart \huge\textbf{H}\normalsize ÆC TIBI scribo fili timotheeisperans me venire ad te cito. P Superfluit fili timothee. Ne putarci timotheus epistolam hanc signum esse quod paulus non efset ad eum venturus, hanc opinionem excludit. SI AVTEM tardauero.P Ecce quare scripsi, vt si tardanero. VT SCIAS qumo oporteat te in domo dei conuersari, quae est ecclesia dei viui. P Aduerte prudens lector quod ecclesiam dei appellat domum deitco quod vniuersa ecclesia enim domus est enim familia est vnum hums patremfamilias. oium. n. christianox conio spiritualis mutua estiquemadmodum eorum qui sunt vnius domus seu vnius familiae. COLVMNA et firmamentum, pro. fundamentumiveritatis.) Ecclesia appellat columna et fundamentam veritatis: pro quamro sostentat veritate  \pend
\section*{CAPVT .IIII. }
\marginpar{[ p.143 ]}\pstart elsam, non qualirercumque sed vt in qua fundata est ipsa veritas nam veritas diuinox mysteriom velut ex celsa sustentatur in ecclesia, dum supra omnia asseritunveneratur et colitur. et in ipsa ecclesia ita fundata est vt extra eam non inueniatur. ET manifeste, pro. et confesse. P hoc est omnium confessione. MAGNVM est pietaris sacramenrum, pro. mysterium.P De ecclesia dicit confessum esse quod est mysterium pietaris. Pietas proprie officium cultumque exhibet parentibus et sanguine iunctis. et quia ecclesiam domum dei dixerat, apte et proprie subiungit quod est magnum mysterium pietatis. Officia siquidem ecclefiae sunt officia pietatis: vrpore domesticortur in domo. Et quia in domo dei tanq̃ patrisfamilias officia ecclesiae sunt officia pietaris erga deum vt patremfamilias. Et quoniam hoc in congregatione fidelium est intimum, ideo mysterium appellatur. Et quia tum extensiua tum intensina quantitare praestat, ideo magnum dicit. P Habes ergo ecclesiam dei viuentis (hoc est congregationem fidelium in deo viuente) relariue ad deum esse domum dei: relatiue ad veritatemjesse columnam et fundamentsi veritatis:relariue vero ad pietatem esse magnum mysterium pietaris, et hoc confesse. vt per vltimum intelligamus officium christi fidelium consistere in pietate, exhibendo parrifamilias cultum et aliis domesticis officium. Vere siquidem talis pietatis eccle sia est mysterium magnum. QVOD manifestatum est.P Vniuersa haec litera vsque ad finem huins capituli nescitur qumo corruptasit. Tum adimendo substantiusi. nam in textu graeco non habetur relariuum quodi sed habetur substantiuum deus. Tum consequenter mutando genus. nam oia quae in neutro genere leguntur (manifestatumjiustificatumipraedicatumj creditumjassumptum) masculini sunt generis, sicut et deus.) Le gendum itaque est. DEVS manifestatus est in carneiustificatus est in spiritul visus est angelispraedicatus est genribus creditus est in mundojassumptus est in gloria.P Et dicunt haec a paulo tum ad manifestandam illam particulam domo deiiet parriculam et confesse:vt confessione oium istorum pateat qd̃ dixerat et confesse, Tum ad explicandum deum patremfamilias erga quem est pietas ecclesiae, dixerat n.quod ecclesia est magnum mysterium pietatis. V Manifestat autem modo ipsum deum sub romne patrisfamilias. dixerat.n.in domo dei.Et vere hic pr̃familias describit dicendo. deus manifestatus est in carne, tum quum verbum caro factum estjet co ruscauit tot miraculis viuens in came mortali:tum in resurrectione manifestatus est in carne imortali. Pm vtrumque namque starum deus manifestatus est in carne. Et hoc est primum qd̃ colit pieras ecclesiae in domo dei pa triffamilias: quoniam per hoc deus factus est paterfamilias per hoc factus est eiusdem naturae nobiscum conmunicans nobiscum etiam in carne. IVSTIFICATVS est in spiritu.P Si miraris deum iustificatum esse, recole dictum esse deo anteq̃ manifestaretur in carneliustificeris in sermonibus tuis. Multo ergo magis poltq̃ manifestatus est in carne, dici de eo potest iustificatus est in spiritu. Et est sensus q in spiritu sancto quem dedit discipulis et consequenter alnsjiustificatus est a calummiis iudaeorum et aliorum qui crucifixerunt a persecuti sunt eum. Vnde in die pentecostes apud ipsosmet qui crucifixerant eum linguis clamando erucifigeincoepit iustificari:vt patet, in actis apostoloz. Et nihilominus adhuc mortalis iustificatus est in spiritu vitaejoperum ac doctrinae ab obiectionibus iudaeorum. VISVS est angelis.P Dens in carne et similiter dens in spirirujab hominibus credi potuit:non autem visus est hominibus. ab angelis autem beatis visus est. Et inseruit hocjad afferendum restimonium angelicum de nostro patrefamilias. Et hoc verificar. tum est de deo in vtroque statujcarnis passibilis et impassibilis. PRÆDICATVS est gentibus.) Ecce do mus dei compraehendens gentes. CREDITVS est in mundo. P Eccelatitudo domus deiicomprçhendens vniuersum mundum. ASSVMPTVS est in glorab vt manifeste intelligamus patremfamilias domus ecclesiasticae regnare gloriosum in gloria ad quam assumptus est in die ascensionis. Haec sunt bencficia patrisfamilias erga quem est pietas ecclesiae. Et simul perpende quod dixerat et confesse, explicatuma carneja spiritujab angelis a gentibus a mundolagloria coelesti.  \pend
\endnumbering\beginnumbering\section{CAPVT .IIII.}\pstart \huge\textbf{S}\normalsize PIRITVS autem manifeste, pro, loquibiliterdicit. P Tractat illam particulamjde fidenon ficta. et dicit quod spiritus sanctus non obscure sed loquendo vt possit intelligi dicit hçesubiuncta. QVIA. pro. quin nouissimis, pro.in posterioribus temporibus discedent quidam a fide. P Praemonemur, vt iacula praeuisa minus feriant. ATTENDENTES spiritibus erroris. pro.erroneis.) Spiritus christi est spiritus veritatis: spiritus autem isti sunt spiritus a vero errantes. quales sunt omnes spiritus haereticorum. omnium enim haereticorum spiritus a vero declinant et errant hac atque illac. ET doctrinis daemoniox. P Spiritus haereticorum pingit ab errore, doctrinas vero a daemoniis: ad significandum quod doctrinae talium hominum non tam sunt doctrinae hominum q̃ daemoniorum. ipsos namque homines appellat daemonia:eo quod officio sunt daemoniajeo q discipuli sunt et imitatores daemoniorti. Et vt clare percipias quod huiusmodi haereticos nominat daemonia, subinngit de eisdem. IN HYPOCRISI loquentium mendacium. pro. in si mulatione falsiloquorum.» Vbi duas conditiones eorundem iungit: et quod falsum loquuntur et quod non ignoranter sed simulate. ET cauteriatam habentium suam conscientiam. P Quoniam metaphoricus est sermo, obstinata horum conscientia intelligi potest ad fimilitudinem indelebilis notae quae ex cauterio facta est in hominibus seu iumentis. vt fignificetur conscientiam hox cauterio diaboli penerratam esse, vt indelebiliter eandem rerineant conscientiamjvt nulla fir spes emendarionis eorum. Est enim sermo de haeresiarchis de capitibus errorum. PROHIBENTIVM nubere. Depictis magistris  \pend
\section*{PRIORIS }
\section*{AD TIMOTHEVM }\pstart erorum describuntur eorundem praecaepta circa nuptias et cibos. Et quo ad nuptias dicitur phibentium nubere quo ad cibos vero. IVBENTIVM abstinere a cibis.P Vtrunque praeceptum dicit esse manichaeotum. OVos deus creauit.b Contra praeceptum abstinentiae a cibisjaffert paulus rationem ex instirurione diuina. AD PERCIPIENDVM cum gratiarui actione fidelibus et iis qui cognouerunt veritatem quia pro. quodjomnis creatura dei bona est: et nihil reiiciendum quod cum gratiarumactione percipitur.) Ne quaereres vbi habetur quod deus creauit omnia ad percipiendum cum gratiarumactione:declarat paulus q in hoc ipso quod omnis creatura dei munda est fm sq et quod nullius vsus illicitus est si cum gratiarumactione fit, ha betur quod deus creauit omnia comestibilia ad percipiendum cum gratiarumactione. Et primum quidem (scilicet quod omnis creatura dei bona est) tanq̃ manifestum relinquit. alterum vero spectans ad vsumiprobat reddendo rationem. SANCTIFICATVR enim per veibum dei et orationem. pro. intercessionem.) Propter iudaeos videtur hoc dicererapud quos cibi aliqui erant imundi quo ad vsum. Aduersus sigdem hoc dicit quod mundificatur vsus ciborum qui reputantur imundi, per verbum dei quo vtimur benedicendo ipsoss et per intercessionem. Radix vnde verbum dei mundificatisignificatur intercessio (hoc est mediario) christi inter legem et non discementes inter cibos. P Er extendit ad legalem prohibirionem sermonem suum paulus:eo quod conueniunt quo ad hoc in materia praecaepriihaeretici et iudaei sectatores legis:dum vtrique im bent abstinere a cibis quibusdam HEC proponens fratribus. pro.haec admonens fratres, bonus eris minister christi iesu enutritus, pro enutriens verbis fidei et bonae doctrinae quam assecutus es.P Vndig con mendat officium admonendi fratres de praedictis, tum ex parte officii relati ad christum. tum ex parte osficii relati ad rem dictam et auditores: manifestando nutritionem animarum in fide et bona doctrina. Quod aurem dicitur quam assecutus esjnon plene explicat textum graecum:significantem quam hucusque securus es. INEPTAS. pro. prophanas autem et aniles fabulas deuita.P Fabellas aduersus religionemjaniles appellar: quoniam similes sunt fabellis quas anus narrant pueris.  \pend\pstart \huge\textbf{E}\normalsize XERCE autem teipsum ad pietatem. P Tractat illam particulamjconscientia bona. Et inchoat ossi cium conscientiae bonae a pietate: quae cultum ac officium exhibet parentibus sanguine iunctis pa triae eiusque beneuolis. NAM corporalis exercitatio ad modicum vrilis est.P Non vituperat corporale exercitium, sed praeponit illi pietatem. Et haec intellige comparando haec fm ppria genera. nam secus est si fm charitatem imperantem aut fm occurrentem articulum necessitatis comparatio fieret. Ex maiori nanq charitare si proueniret corporalis exercitatio, vtilior esset q̃ pietas ex minore charirate. et similiter in arti culo necessitatis vrgentis ad ieiunium seu aliquid huiusmodi corporalel vtilior esset corporalis exercitario q̃ pietas. Sed haec sunt per accidens, nam fim genus pprium corporalis exercitatio in ieruniis et absti nentiis peregrinationibus ciliciis et huiusmodi, ad modicum vtilis est:quia tantum valet ad domandum cor pus proprium. PIETAS autem ad omnia vtilis est. P Officia pietatis erga ecclesiae membraj et cultus erpa patremfamilias domus quae est ecclesia, sunt vtilia ad omniajtam temporalia q̃ spiritualia. Vnde hoc manifestat subiungendo. PROMISSIONEM. pro. pmissiones habens vitae quae nunc est et future.P Ex hoc ipso q in lege diuina in praecaepto pietaris (honora patrem tuum et matrem tuam )exo.2o. adiungitur promissio vitae praesentis et futurae, manifestat paulus quod pietas ex suo genere ad omnia vrilis est. Interpretal autem paulus verba legis partim ad literam et partim mystice. et propterea pluraliter dicit promissiones et vtriusque vitae meminit. In exodo namque dicitur vt sis longçuus sup terram quam dominus deus tuus dabit tibi. paulus vero etiam mystice terram promissionis coelestem patriam interpretarur. et proprerea dicir virae prae sentis et futurae. ¶ Nec hinc arguas inferendo abstinentiam et ieiunia inutilia esse ad vitam aerernamtqum hoc hinc non habetur:sicut nec habetur reliqua virturs̃ opera inurilia esse ad viram aeternam. Tantum hinc habetur speciale priuilegium pietatis;quod in illius praecaepto explicatur promissio vtriusque virae, cum hoc enim priuilegio stat et reliquarum virtutum opera (inter quae est exercitatio abstinentiae corporalis) meritoria esse virae aetemae.P At si quaeris cur pietas donata est hoc priuilegio. rario est quia pieras ad faciendum alteri bonum tendit, quod diuini officii est maxime imitatiuum et homini difficile: quum naru raliter non inclinetur nisi ad proprium bonum. V FIDELIS sermo et omni accaeptione dignus.P Ad declarandum quod ipse interpretatus fuerat (videlicet et futurae) adiungit quod sermo iste dicens quod pietas hubet promissionem vitae futuraeest sermo fidelis et omni accaeptione dignus. Et hoc ab experientia in seipsis manifestat, subiungendo. IN HOC enim.P Deficit. ETi laboramus et maledicimur.p hoc est nam ad hoc pietatis opus et laboramus assidue et maledictis afficimur. QVIA speramus. pro.sperauimus in deum viuum. pro.in deo viuentequi est saluatoromnium hominsi, maxime fidelium. F Ab experientia qua sustinent rum proprios labores tum maledicta exercendo ea quae pietatis sunt, manifestat promisi sionem vitae futurae:ex eo quod haec omnia pcedunt ex spe in deo viuente quod saluaturus est omnes homines praecipue fideles. In vita futura saluabit deus omnes homines quo ad vitam corporis suscitando omnes. maxime aurem saluabit fideles: quoniam suscitabit eos ad vitam foelicem. Haec est res sperata:de qua dicit. sperauimus in deo viuente qui est saluator etc. pro qua et laboramus et maledicimurjvacando pietari eccle siastice. CPRÆCIPE haec et doce.) doce vt sciant, praecipe vt imitentur. NEMO adolescentiam pro. iuuentutem tuam conremnat hoc est sic te habe vt nemo contemnat aetatem inuentutis tuae. Vnde et  \pend
\section*{CAPVT }
\section*{V }
\marginpar{[ p.144 ]}\pstart statim aduersatiue subiungit. SED exemplum esto fidelium in verbo.] hoc est in loquendo. IN con uersationejin charitate.* Praecipit paulus vt locutio sit exemplaris, conuersatio sit exemplarisjdilectio sir tremplaris. Et praeter haec adiungit tria alia. Quox primum deficit in textu latino: videlicet. IN SPIRI T vjin fideuin castitate. P In spiritualibusjin iis quaesunt fidei et in corporis castitare, praecipit etiam timotheo vt exhibeat se exemplum fidelium. Et hae sunt partes ex quibus effici debet vt nemo eius iuuentutem con temnat. Et multo magis hae partes in praelato puectae çtatis efficient vt non despiciatur.  \pend\pstart CDVM venio attende lectioni.] ad pascendum teipsum. EXHORTATIONI. P aliorum ad bona ope ra. ET doctrinç.p Superfluit et. non solum exhortare sed doce alios ea quae sunt sacrarum scripturarum. NOLI negligere gratiam quae. pro. donum quodiin te est:quae data pro. quod datumjest tibi per prophe riam cum imposirione manuum presbyteni. P Episcopalem gradum appellar donum datum timorheo duplici concurrente causa. Altera est prophetia quod futurus esset episcopus. a quo autem fuerit facta haec prophetia seu cui fuerit reuelatasscriptum non est. Simile quid in.1.ca. dixeratifm praecedentes ad te prophetias. PAltera est impositio manuum presbyterii:hoc est officii presbyteratus. ad fignificandum impositionem non esse humanae manns, sed manus sacraejmanus sacerdotalis officii. Et dicit pluraliter manuum presbyterufforte ad significandum plurium sacerdotum concursum. fiebant siquidem episcopi per impositionem manuum, sicut et hodie fiunt.V HÆC medirare. pro. exerce. P Dixerat praeceptum negatiuumjne neglexeris:adiungir modo affirmatiuumjhaec exerce. Et pronomen haęclmonstrat tria antedicta episcopi officialegerejexhortari et docere. IN HIS esto vt profectus tuus manifestus sit omnibus pro.in oibus.P Fructus consequens suapte natura explicatur. Nec hic fructus vacat salute aliorum VATTENDE tibiipsi et do ctrinae. P Adiungit vt sic attendat saluti aliog vt simul incumbat propriae saluti. INSTA. pro. persister in illis.P vtrisque :scilicet quae sunt ad salutem tuam et aliorum. HOC ENIM faciens.P Attentionem ad vtraque demonstrat pronomen hoc. ET TEIPSVM saluum facies et eos qui te audiunt.P Ecce fructus attentionis ad vtrunque : et vere fructus perfectus. \pend
\endnumbering \beginnumbering\section{CAPVT .V.}\pstart \huge\textbf{S}\normalsize ENIOREM ne increpaueris, sed obsecra. pro. exhortarejvt patrem.P Nomen aetatis esse hoc in loco senioremjex aliis aetatibus subiunctis apparer. Et propter aliorum aetates subiunctasjmagis quadrare viderur exhortare q̃ obsecra: quoniam ad omnes subsequentes idem verbum referf.  Peccantem itaque senemimandat paulus non increpandum sed exhortandum: vt veluti filiali vtatur in eum reuerentia.sic enim facilius corrigitur senex. IVVENES vt fratres anus vt matressiuuenculas vt sorores. P Repete exhortare. IN OMNI castitate. P hoc opus fuit addere in exhortatione iuuencnlarum propter periculum aetatis. Et dicit in omni, vt sit castitas in oculisjin gestibusjin verbis etc. VVIDVAS honora.P Honorem in scripturis sacris referri non solum ad officia sed etiam ad subsidiajdo cet dums exponendo praeceptum honora parentes tuos de subuentione. Nec ambigendum est vtrumque hic comprehendi. verum sub ratione honoris praeceptum datur:quoniam honor per se exhibendus est, subuentio aurem per accidens: videlicet si sunt indigentes. V Viduis autem honorem exhibendum iubet non omnibus, sed QVE vere viduae sunt.P Et est aduerbium vere tam hic q̃ in sequentibus similibus. Quas autem vere intelligat viduas ipsemet explicabit ex duplici differentia altera penes naturalia. altera penes volumtaria. Vnde a naturali necessitudine inchoans dicit. SI QVA autem vidua filios aut nepotes habet.P Prima differentia inter viduas tangitur:consistens in hoc quod aliqua vidua est quae habet vinculum natiiralis necessitudinis cum filiis aut nepotibus (hoc est filuis filiorum) aliqua autem caret huiusmodi vinculo. Et de habente huiusmodi vinculum subiungit. DIS CAT. pro. discantiprimum domum suam regere, pro. pie curarej et mutuam vicem reddere parentibus. P Dicendo discantjmanifeste loquitur de filiis et nepo tibus. Dicendo primumssignificat quod vidua non prius recipiatur ad professionem vidualem q̃ filii sen nepotes didicerint curam habere domus et mutuam vicem parentibus reddere. Et in promptu est ratio huius praecepti: quoniam quandiu filii seu nepotes non peruenerint ad aetatem et virtutem regendae domns, incumbit viduae cura familiae: et hinc impeditur ad vacandum diuinis et ecclesiae, quibus vacabant pfessae viduitarem. Tales enim viduae q̃uis forte affectu sint vere viduç:quia tamen alligatae sunt curae filiorum non sunt recipiende ad professionem vidualemjdonec filii sen nepotes discant domum suam regere et mutuam vicem reddere parentibus, habendo curam parentum qualem parentes habuerunt de eis:hoc est donec filii ha beant curam viduae, vt sic vidua exempta ab omni cura domestica possit profiteri statum vidualem ad vacandum diuinis hospitibus et aliis misericordiae opibus. HOC enim. P Deficiunt duae dictiones. Bo NVM etjacceptum est coram deoper Ne vidua talis aegre ferat dilationem professionis, explicat paulus hu ausmodi dilationem bonam esse in se et gratam deo: quoniam innititur necessariae pietati.P QVE autem vere vidua est. b Officia eius quae vere est vidua describit, tum ad declarandum quod dixerat vere viduae sunt. tum ad explicandum incompossibilitatem viduae habentis curam nepotum ac filiorum ad officia vere viduae. ET DESOLATA speret. pro. speranirjin deum. pro.in deo, et instet, pro. et persistirun obsecrarionibus et orationibus nocte et die.) Ecce duo officia vere viduae. tum quod desolata locauit spem suam in deo, non in futuro coniuge non in diuinis et reliquis quae mundi funt tum quod tempus omne occupat oratio-\pend
\section*{PRIORIS AD TIMOTHEVM }
\pstart nibus et obsecrationibus. quod non porest facere vidua alligata curae filiorum minorum. NAM QVE in delitiis est, viuens mortua est. P Alteram tangit differentiam viduarum penes voluntaria. E regione siquidem ad persistentem die ac nocte in orationibus describit eam quae delitiis vacat. Et de hac dicit q viuens mortua est: hoc est vtrique subiacet contrario.nam et viuit statu, mortua est aurem opere. Ira que quatenns viduaimonstruosa est nam et viuit et morrua est. viuit nomine et statu viduae:mortua est carens ope ribus viduarum.V ET HOC praecipe, vt irrepraehensibiles sint. P Ne parnifieret huiusmodi mors et vita, subiungit hoc spectare ad repraehensibilitatem viduarum. Et propterea mandat timotheo vt praecipiar vi duas abstinere a delitiisjvt irrepraehensibiles sint. SI QVIS autem suorum et maxime domesticorum curam non habet. pro.si quis autem suis et maxime domesticis non prouider, fidem negauit et est infideli deterior.) Relariuum quis q̃uis graece anceps sit ad masculinum et foemininumj et propterea legi possit si qua referendo viduam (nam de viduis est sermo) et sic legendum crediderim (tum propter subiuncta de officiis viduarum. tum propter praecedentia de officio viduae habentis filios de quo quum dixisset hoc enim bor num et accaẽptum est coram deo, ne putaretur illud sic esse bonum gratumque deo quod etiam pfiteri viduitatem esset bonumjexplicat modo quod esset magnum facinus) fateor tamen quod si relatiuum communis generis ha beremus magis quadraret: vt ex vniuersali canone communi vtrique sexui monstraretur quantum crimem est non prouidere filiis et nepotibus. In cuius signum non loquitur canon de filiis aut nepotibus sed vni uersalissime de suisjpraeponendo domesticos. Y Quod autem dicit fidem negauitjad exaggerarionem sceleris dicitur. et verificat quia fidem negauit facto. Quod vero subiungit et est infideli deteriorjexpomne non indiget: quia manifestum est. nam infideles puident suisjhoc est spectantibus ad se, et maxime domesticis. ¶VIDVA eligatur non minus sexaginta annorum pro Dixerat viduas honora quae vere viduae sunt: et descripserat eas quae vere sunt viduae et quae mortuae viduae et quae impeditae ad tempus declarat mon quae inter vere viduas praehonoranda sitiquae eligenda sit in caput viduarum. Et multas exigit conditiones. Quarum prima est aeratis. Et haec tantae aeraris exactio manifestat quod de electione capiris est sermo. QVE fue rit vnins viri vxor.3 Castitatis muliebris decor est post primum virum non nupsisse. IN OPERIBVS bonis testimonium habens. P Non dicit testes sed testimonium sui habens in ipsis operibus bonis. Et quae sint illa bona opera explicarssubiungendo. SI FILIOS educauit.p Intellige de ea quae habuit fi lios. Si enim filios neglexisset!malum haberetur de ea testimonium. SI HOSPITIO recaepit.P Hospitalitas magnum est restimonium charitatis. SI SANCTORVM pedes lanit.P Pedires ibant praedica tores et enangelici viri (quos paulus sanctosjhoc est mundos appellat) quor pedes lauare tunc consueue rant viduae, aut recipientes eos hospitio:aut si alibi hospitabanturiilluc veniebant ad lauandum corum pedes. SI TRIBVLATIONEM patientibus subministrauit. P Et hoc est alind misericordiae opus subuenire afflictis. Et demum SI OMNE opus bonum subsecuta. pro. secutalest. per Ne persingula moraret vniuersaliter subdit si omne opus bonum est secuta, faciendojprocurando etc. Et in his officiis perspicert potes quae sunt officia vere viduarum relatiue ad hoies. nam relatiue ad seipsas descripserat abiiciendas de litias: et relarine ad deumsque persistit in orationibus et obsecrationibus die ac nocte.  \pend\pstart ADOLES CENTIORES. pro. iunioresjasĩt viduas deuita.P Post impedimentum temporale (videlicet q vidua habet filios aut nepores educandos) impedimentum aliud mains edocet:aetatis. s.nuenilis. Quam vitandam dicit quo ad professionem vidualis status. Non dicit differendam esse professionem (sicur direrat de habente filios) sed absolute dicit deuita. Er rationem subdit. QVVM enim luxuriatae fuerint in chio. Superfluit enim et iuxra textum graecum geniriui est casus christo. Et sumitur luxuriari propriejpro supfluo vsu. Et dicit quum enim luxuriatae fuerint christi, non quod superfluitates vere fuerint christi, sed quod vsae sunt christo in superfluitatibus. Significatur enim ad literam quod quum in superfluitatibus vsae fuerint viduali velo christi, tunc NVBERE volunt.P Haec est itaque vera rario repellendi iuuenculas viduas a profelsione viduali:quia quum vsae fuerint aliquanto tempore velo christi in superfluitatibus mutant voluntatem et volunt nubere. HABENTES damnationem, pro. indicium quia primam fidem irritam fecerunt. pro. abie cerunt. P Tria proposuerat. s. deuitajquum luxuriatae fuerintj et volunt nubere. Et retrogrado ordine illa ex planat: dicendo quod nubendo incurrunt iudicium iustae damnationis eo quod repulerunt primam fidem. Manu festant autem subiuncta de eisdem (volo igitur iuniores nubere) quod non est sermo de fide primi coniugis absolute, sed de eadem firmata professione viduitatis haec est enim prima fides cuius abiectione incurrunt iudicium. Et hinc sumitur quod profitebant viduiratem voto in manu episcopi. et a tali voto vitandas ado lescentiores viduas paulus censet propter periculum instabilitatis: quia postea volunt nubere. SIMVL autem et ociosae discunt circumire domos. P Declarat luxuriatae fuerint christi, explicando in quibus con sistit hmoni superfluitas vtendo christo in huiusmodi superfluiraribus. Et primum suppurat ociole dircumire domos. Coniicitur hinc quod officium vere viduarum quae erant professae viduitatemjerat ire per domos, non ociose sed pietatis causa. et proprerea dicit paulus quod harum exemplo istae adolescentiores viduae discunt circumire domos, non tamen pie sed ociose. Et hic est primus luxus. NON solum. Deest, autemjociosae sed et verbose. pro. nugaces! et cunosɇ, loquentes quae non oportet.P Ecce multitudo luxusi in nugis et curior fitate, et (quod pro eius est in verbis indecentibus hoc est enim vti christo in superfluitatibus. ¶ VOLO autem  \pend
\section*{CAPVT }
\section*{.V. }
\marginpar{[ p.143 ]}\pstart pro. igitur iuniores nubere.) Declarat illam particulam deuita: inferendo ad euitandum tum periculum damnationis earum tum huiusmodi luxus quod non admittantur ad professionem viduitatis sed nubant.Et hinc vt dictum estjmanifeste patet quod licite nubebant ante professionem viduitatis. FILIOS proereare. Et hinc manifeste patet quod de iunionibus vere iuuenculis loquitur, non de iunioribus. i. minoribus sexagenanis. nam de illis nonfuisset dictum filios procreare. MATRESFAMILIAS esse, pro. familiam admi nistrare.b Et hoc contra ocium et discursum. NVLLAM occasionem dare aduersario maledicti gratia. b Dabatur occasio maledicendiimurmurandi ac detrahendi, ex discursu iuueneularum viduarum perdomos. et propterea ad hoc cuitandum dicit quod tales iuueneulae nullam occasionem dent aduersario(sine gentili siue christiano aduersanti) maledicentiae. IAM enim quaedam conuersae sunt retro.P Loco harum trium dictionum (conuersae sunt retro) legendum est. DEFLVXERVNTI post satanam. P Ecce effectus manifestans non solum periculum sed casum quarundam talinmjet simul occasionem datam maledicntiae. Et ex hoc rertu apparet quod ante hanc pauli inhibitionem admittebantur iuniores viduae ad professionem vidualem, et sequebantur huiusmodi inconuenientia praedicta:quibus cognitis paulus inhibuit huiusmodi innem culas viduas admitti ad professionem viduitatis. Didicit ab experientia paulus non conducere illis iuuenculis viduis nec conducere ecelesiae huiusmodi professionem. Vtinam hodie disceremns ab huiusmodi experientiis an prosint iuuenilis aetatis vtriusque sexus personarum solennia votajtum sacrorum ordinum tum religionum.  \pend \pstart SI QVIS fidelis. P Desunt hic quattuor dictiones. AVT si qua fidelis habet viduas, subministret illis vt, pro. et non grauetur ecclesia:vt iis quae vere vidnae sunt sufficiat. P Hactenus de spiritualibus circa vi duas locutus fuerat in speciejmulta instruendo:modo de temporalibus subuentionibus instruit distinguem do inter viduas vere desolatas et vere viduas non desolatas. Et non desolatas appellat illas quae habent suos iam christianos: vere autem viduas (hoc est desolatas) appellat quarum sui aut non suntjaut non sunt christiant Et de illis quidem mandar vt sui fideles subministrent illis temporalia, et non grauerur ec clesia in communi pro illis: vt ecclesia sufficiat ad subueniendum iis quae vere desolatae sunt. Perpende hinc σ absque delectu ecclesia in communi proudebar viduis indigentibus. et quoniam hoc onerosum erat eccle siae, ordinat paulus vt filiiinepotes cognati et huiusmodi prouideant suis viduis: et curae ecclesiae relinquantur solae vere desolatae. Et ratio mandati huius est vt suauiter continuari possit ecclesiae communis cura circa viduas.  \pend\pstart \huge\textbf{Q}\normalsize VI BENE praesunt presbyteri. pro Nomen ordinis est presbyteri hoc in loco. Quod ex eo patet q describuntur praeesse bene praeesse enim est episcoporum feu superintendentium.  Nec dicit qui praesunt sed qui bene praesunt. DVPLICI honore digni habeantur.p videlicet reuerentialis officii et subuentionis.hic est enim duplex honor in scripturis. MAXIME qui laborant in verbo et doctrina.b Inrer bene praesidenres praeferendi sunt illi qui super officium superintendendi adiungunt laborem in ver bo exhortarorio affectus, et in doctrina instruenre intellectum. Et ad hoc affert autoritatem scripturae deurero.25. DIXIT. pro. dicitjenim scriptura. non alligabis os boui trituranti. P Nora trituranti:quoniam propter hanc particulam affertur autoritas ad eos qui laborant in verbo et doctrina. ET DIGNVS est operarius mercede sua. P Hoc scriptum est in euangelio lucae. 1o. Et ad idem affertur operariusinon ociosus ADVERSVS prelbyrerci. P Consequenter presbyteri nomen non aetaris sed ordinis apparet:quoniam de talibus est praesens tractatus. ACCVS ATIONE M noli recipere, pro.non admittejnisi sub duobus aut tribus testibus.b Non facile admittendam censet accusationem presbyteri ratione status. non enim vult vt accusarione admissa subiaceat discrimini progressus an sint testes sufficientes an deficiant reftes, sed quod ipsa accusatio non admittatur nisi constet de sufficientibus et idoneis testibus. Et hoc est pri uilegium presbyterorum exhac pauli ordinatione. PECCANTES coram omnibus, arguervt et caeteri timorem habeant. P Anceps est constructio an coram omnibus referatur ad peccantes an ad argue. Et si quidem ad peccantes refertur, de punitione peccatorum publicorum instruit timotheum vt peccantes coram omnibus (quorum non oportet accusationem et testes expectare) arguatjpuniendo. Equadrat ratiojvt et caeteri timorem habeant. Si autem refertur ad argue, quum dixisset de presbyteris q accusationem aduersus eos non admittat nisi sub duobus aut tribus testibus, consequenter subiungit quod eos conuictos in peccatojpublice puniat vt caeteri timorem habeant. V Verum si participium peccantes ponderatum fuerit, sola prima lectio quadrare videtur. non enim dicit eos qui peccauerunt sed peccantes. et ideo adiungit coram omnibus, ad differentiam peccantis in teipeccantis secreto: quem do minus corrigendum inter te et ipsum praecipit. CTES TOR coram deo et. Deest dominojchristo iesu et electis angelis vt haeccustodias. ] Cerne quanto vinculo astringit timotheum in materia iudicii. SINE praeiudicio.y hoc est sine anteiudiciossine praeuio indicio. Praecipitationem, inconsultationem indiscussionemque prohibet praecipit discerni prius examinari et discuti causam anteq̃ procedatur ad fententiam. NIHIL FACIENS in alteram partem declinando. P Loco istarum quattuor dictionum (in alteram partem declinando) legendum est. PER ADVOCATIONEM. P Exparte sui errat si procedit absque prae iudicio ex parte autem aliorum errat si aliquid facit per ad uocactionem :hoc est ad in-\pend
\section*{PRIORIS }
\section*{AD TIMOTHEVM }\pstart tantiam aduocatilad instantiam cuiuscunque fauoris patrocinii et reliquorum huiusmodi. Haee sunt propter quae tanto vinculo astringit timotheui in iudiciis. VMANVS cito nemini imposuens. Cir ca promotionem diaconox et presbyterorum instruit ipsum. fiebat enim per impositionem manuum. Et rationem subdit. NEQVE communicaueris peccatis alienis.P peccata siquidem illius communia enunt tibi si minus idoncum pmoueris. TEIPSVM castum custodi.P fic curam gere aliorum vt non recedas acu stodia tuiipsius. NOLI adhuc aquam bibere.P Non solum animae sed corporis timothei curam erercae paulus. SED modico vino vtere propter stomachsi tuum et frequemres tuas infirmitates.5 Snia est clara CEQVORVNDAM hominum peccara manifesta. pro. praemanifestajsunt praecedentia ad iudicium: quosdam autem et subsequuntur. P Dixerat vt nihil faceret in iudiciis sine praeuio iudicio:ne haec intelliger rentur vniformitersed fm diuersitatem materiae subiectae diuersimode seruanda, distinguit peceata hominum iudicanda dupliciter inueniri. vel ante iudicium manifestaivel quorum manifestatio sequir indicium. Et est sermo de iudicio discussionis:sine quo praecaeperat nihil fieri, vt ex hoc dinerso modo mani festationis discerneret timotheus diuersum modum praeuii iudicii. nam in iis quae anteq̃ discuriant ma nifesta sunt, sola notitia apud iudicem opus est:in aliis autem notitia iudicis pendet ex discussione praeuia Est itam ordo literae. Quorundam hominum peceara praecedentia ad indicium hoc est quae indicanda lunt, praemanifesta suntjante scilicet q̃ veniant in iudicium: quosdam autemniudicandos, et subsequuntunscil cet peccata quatenus manifesta. Ita quod peccata quatenus manifesta subsequuntur iudicium discussionis SIMILITER et facta bona manifesta. pro. pmanifesta sunt: et quae aliter se habent.P Superfluit se. ABr SCONDI non pofsunt. P Dixerat nemini cito manum imposueris: distinguit modo de bonis factis su cut distinxit de peccatis. Et dicit quorundam facta bona praemanifesta suntjanteq̃ homines discuriantun vt intelligat timotheus quod in iis qui sunt notorie boni quo ad factajnon est opus discussione ad hoc v imponantur eis manus. Quidam autem sunt qui habent aliter bona:hoc est non notorie bona habent opera. et hi discutiendi suntiquoniam latere non possuntisi diligens adhibeatur discussio. Y Et intellige ipsos latere non posfe regulariter loquendo, communiter enim diligenti discussione inuenitur veritas talium, contingere tamen potest praeterregulariter quod veritas occultaretur. Er ppterea circa hos promonen dos vtendum est tarditate, vt veritas inueniatur.  \pend
\endnumbering\beginnumbering\section{CAPVT VI.}\pstart QVICVNQVE sunt sub iugo seruip De corporalicer seruis loquit. DOMINOS suosomm honore dignos arbitrentur. P. Et loquitur de dominis etiam paganis baptisma enim non absoluit seruos a seruitutis corporalis vinculo. Vnde et rationem subdit. NE NOMEN domini pro. deild doctrina blasphemetur.P Si serui facti christiani non honorarent dominos suos, periculo blasphemiae exponeretur tam nomen dei q̃ doctnna dei, conquererentur siquidem domini de tali doctrina quae subducit seruos a debito seruitio dominom:et similiter detraherent tali deo cui placent huiusmodi iniustitiae. Quo contra si serui christiani omni honore dignos existimant dominos suos, occa sionem dant landandi nomen et doctrinam dei christianorum.! QVI autem fideles habent dominos, non contemnant. P supple ipsos dominos. QVIA fratres sunt. P in christo. SED magis seruiant quia si deles sunt et dilecti quia pro. quil beneficii participes sunt.P Ordo literae est. sed magis seruiant qui be neficii participes sunt, quia fideles et dilecti sunt eorum domini. Et est sensus quod fraternitas in christo nonim ducar contemptum dominorum tanq̃ parium: sed serui qui simul cum dominis participes sunt beneficii diuini, magis seruiant dominisjeo quod non solum domini sed fideles sunt et dilecti a deo. Et meminit participationis beneficii, vt intelligant serui fraternitate non esse subtractum dominiumjsed constitutos esse sen uos et dominos participes diuini beneficii spiritualis vitae. et intelligant augeri debere seruiria exhibenda dominis fidelibus eoipso quod domini aucti sunt: dum ex dominis effecti sunt domini fideles et dilecti a deo. HÆC docet et exhortare. p doce nescienres: scientes vero exhortare ad exequendum. SI QVIS aliter docet.) tanq in fauorem libertatis. ET NON acquiescit, pro, accedirisanis sermonibus domini nostn iesu christi.p Et non accedit docendo. explicat enim quid est aliter docere. Et appellat sermones domini sermones sanos, quia vere sunt sani non solum per elongationem ab omni insaniassed per affirmationem verae sanitatis mentis. ET EI quae fm pietatem est doctrinae.P Sermonibus christi adiungit doctnni fm pietatem, ad significandum quod doctrina verarum virtutum comes est sermonum christitvt intelligamus quod seruos parifacere dominos suos et aduersatur sermonibus christi et doctrinae pietatis. Sub pietate si quidem compraehenduntur in scriptura omnia officia debita proximis: vt ex eo patet quod vnicum tantum praeceptum officiorum debitorum proximo datum est in secunda tabula sub praecipuo pietatis actu qui elt honorare parentes. Loquitur enim ad literam textus hic de pietate quae virtus est exhibens cultum a omn cium parentibus, patriaeisanguine iunctis parriaeque beneuolis: extensa tamen vt dictum est ad omnia offi cia debita proximis. SVPERBVS est. pro. superbinit.p si quis aliter docet etem, superbiuitjiam ex praece denti superbia aliter docet. NIHIL sciensP Superbiuit in affectu, nihil sit in inrellectu. SED LANT GVENS circa quaestiones et pugnas veborum) NIhil scit rei veritatis, sed est ad similitudinem aegro\pend
\pstart tantis cirea in quisitiones et pugnas verborum: puta quia christus nos liberauit, habent mentem velut aegrotam ad inquirendam vim vocabuli libertatis ad discernendum pugnam verbog libertatis et seruitutis. tanq̃ non stet simul libertas qua christus nos donauitjcum seruitute ciuili. Ex superbia enim talinm doctor puenit et quod nihil sciunt quo ad ventatem reiiet quod circa inquisitiones et pugnas verborum non habent mentem sanam sed aegrotantentac per hoc deficientem etiam a veritate circa verba. et propterea oino nihil sciunt:qa necres nec verba. EX QVIBVS.P superbiaiquaestionibus et pugnis verborum. ORIVNTVR inuidiae, pro. inuidi ajcontenriones. pro. contentiojblasphemiaeisuspiciones malaejconflictariones. P Vide quot mala inde sequunt. HOMINVM mente corruptorum et qui veritare priuati suntjexistimantium quaestum esse pietatem.] Dixerat quod ex his oriunt tot mala:explicat mon in quibus hoibus ortus horum malorum locum inuenit. Quum mens sit incorruptibilis, mente corruptos appellat non fm substantiam, led fm ignorantiam prauae dispositionis. Vnde effectum huius corruptionis explicat dicendo a veritate destitutox. sunt. n.tales velut deserti a veritateiet rerum et verborum. Et quod in genere dixerar de corruptione mentis et destitutione a veritate, manifestat in indiuiduo quum dicit existimantium quaestum esse pietarem. Nunq̃ christus docuit discipulos suos quaestum, nunq̃ ordinauit vt quaererentinunq̃ vt mendicarent. Paulus mente corruptis et a veritate destitutis attribuit doctrinam dicentem quaestum esse pietatem. quaestus n.est aut necessitatis aut auariciae:nunq̃ autem pietatis. Et est sermo de quaestu pro se:secus autem est de quaestu pro pauperibus in hierusalemP Deficiunt in textu latino tres dictionesistae. ABSISTITO a talibus.P Non solum vitandos sed procul esse a talibus mandat paulus. EST autem quaestus magnus pietas cum sufficientia.) Illi corrumptae mentis hoies existimabant quaestum temporalium rerum esse pium. paulus e regione docet non solum quaestum sed magnum quaestum esse pietatem:qm̃ est quaestus magnox bonorumj est quaestus spiritualium, quae sola sunt vere magna bona.V Adiungit autem pietari sufficientiam ad eradicandum quaestum remporalium: vt intelligamus sufficientiam victus et vestitus non quaeri a pietate, sed concomitari pietatem. Et vere sic dominus ordinauit, vt discipuli pietatem exercentes praedicandoisufficientiam habe rent non quaererent. edentes inquit et bibentes quae apud illos sunt. offerri vult pietatiinon quaen ab illa. NIHIL enim intulimus in hunc mundum.P Superfluit hunc. HAVD dubium quia pro.quodinec auferre ali quid possumus.p Declarat comitem pietaris sufficientiam quae comprçhendat: discurrendo ab ingressu nostro in mundum vsque ad egressum per mediam vitam. Et quia in ingressu nihil facultarum temporalium intulimus, non est dubium quod similiter in egressu nihil auferre possimus. in vita vero media. HABENTES autem alimenta et quibus regamur his contenti simus. pro. erimus. p Ecce quae compraehendit sufficientia:alimenta scilicet et regumenra qumquidem his contenti crimus fm necessitatem naturae Ex hoc.n.quod natura contenta erit in hac vita alimentis et tegumentis, monstrat has esse partes sufficientiae quae est comes pietaris.  \pend
\section*{CAPVT }
\section*{.VI. }\pstart NAM QVI volunt diuites fieri. P Eregione non contentox alimentis et regumentis!locat eos qui volunt ditescere. et eos pingit suis coloribus valde malis. INCIDVNT in tentationem. P auariciaeffraudis! et hmoni. ET IN laqueum diaboli.p Superfluir diaboli et prepositio enim Non solum succumbunt tentationi allicientium diuitiarum ad superflnitatem, sed etiam illaqueant laqueo pene insolubili:eo quod obnoxia restitutioni tales nunq̃ restituunt. sunt. n. velut illaqueati affectu rerum. ET desideria multa inutilia. pro.irrationabilia. per Et vere sic esse apparet: qm̃ in multa incidunt desideria aliena a recta ratione. Er quod peius est. ET nociua.p aliis. Haec.n. coniter videntur accidere iis qui volunt ditescere. Et bene nota ne pistos intelligas eos qui pro necessitate familiae suae industrias exercent mercaturae. hi.n. non sunt qui volunt di tescere sed lucrari vt possint fm statum pprium vinere filii etc. QVÆ demergunt hoies in interitum. P animae in psenti. ET perditionem. p in futuro. RADIX enim oium malog est cupiditas. pro. auaricia.P Dictio graeca non est nomen cupiditatis in genere sed cupiditatis argenti, quam vnico vocabulo auariciam dicimus: qm̃ de inordinata cupiditate est sermo. Et hçe metaphorice dicit radix malorum oium, qm̃ ex ca pullulant mala omnia. Sex.n.genera malox paulus iam supputauit ex tali radice oriri:incidere in tentationem et laqueumjdesideria multa irrationabilia et nociuajinteritum et perditionem. Sub his generibus mala omnia continentur. QVAM quidam appetentes aberrauerunt a fide.P Manifestat ab experientia auariciam radicem esse oium malorum, vsoeadeo vt sit eriam radix infidelitatis. Vehementer siquidem animus affectus di uiriis etiam a fide apostatat. ET inseruerunt se doloribus multis. pro Nihil paulus pretermittit. Declarauit quod auaricia est radix oium malorum culpae ex experientiaiquia abducit etiam a fide:declarat mon quod est etiam radix multog malorum poenae in hac vita, dicendo quod immiscuerunt se doloribus multis. Perturbationes animi quas auari incurruntjconstat esse multas: quoniam aestuant continuo apperitu lucri. Et multo plures perturbationes animi incurrunt tales aberrantes a fide.nam non solum perturbantur propter diuitias mon perditas modo periclitantes etc. sed etiam propter synderesim ob fidem relictam.  \pend
\marginpar{[ p.146 ]}\pstart V AVTEM o homo dei haec fuge. ) Appellat eum hominem dei ad reddendam rationem fugae iuxta T charitatem patientiam mansuetudinem. ) Et cur has virtutes sectari praecipiat subiungit summando ad duo capita. CERTA bonum certamen fidei, appraehende vitam aeternam. ) Ad haec duo reducitur omnis sequela cuiusque virtutis. Et primum quidem spectat ad viam praesentis vitae alterum autem ad terminum eiuschristi praeceptum non potestis deo seruire et mammonae. SECTARE vero iustitiam pietatem fidem  \pend
\section*{PRIORIS AD TIMO. }
\section*{CAPVT VI. }\pstart dem, Sectando siquidem virtutesscertamus certamen fidei bonuĩ: et in termino appraehendimus vitam aetera namjpremium certaminis. IN QVA. pro. in quam. Deest. et vocatus es et confessus es bonam confessionem coram multis testibus. P Ne impossibile reputaretur tam ardnum praemiumldicit. in quamjvitam aeternam, tum vocatus es a deojtum tuipse confessus es confessionem bonam coram multis testibus. Professionem quam facit episcopus supintendendi aliis appellat confessionem bonam factam coram multis testibus: qu non secreto sed palam fit huiusmodi professio. Et hanc quog affert factam in vitam aeternam, Ita quod tam ex vocatione dei q̃ ex pfessione tua certandum est tibi certamen fidei ad apprçhendendam vitam aeternam. VPRÆCIPIO tibi coram deo.P Vinculo vocationis et propriae pfessionis adiungit vinculum praecepti adhibitis testibus deo et christo. QVI viuificat omnia.P Quoniam praecipit media ad vitam aeternam, ideo deum inducit quatenus viuificat omnia. ET CHRISTO iesu qui testimonium reddidit. pro. testa tus estisub pontio pilato bonam confessionem. p Dixerat et confessus es bonam confessionem coram multis testibus. et propterea prçcipiendo executionem illius affert christum iesum quatenus testatus est sub pontio pilato bonam confessionem. Facto siquidem et verbo christus bonam confessionem testarus est sub pontio pilato. Verbo quidem testatus est confessionem propriae innocentiaemon haberes aduersum me porestatem vllam etc. futurae vitaejregnum meum non est de hoc mundo. propriae deitatis lego ad hoc veni in mundum. et demum propriae doctrinaejvt restimonicĩ perhibeam veritari. Facto autemjpariendo vsque ad crucem. Vñ vere de ipso dicit quod confessionem bonam non solum ptulit, sed et testatus esticertando facto passionis suae et verbis testificando illam. VT SERVES. pro. seruare te.P Ecce actus praecaeptus. MANDATVM.Y proculdubio tibi impositum in tua professione. Mandatur siqdem epreo quum ordinat vt praedicet populo dei.ita quod non solum exprofessione propria sed ex mandato obligat. Et ppterea paulus dicitipraecipio tibi ser uare te mandatum. SINE maculasirrepraehensibile, pro. imacularumjimepraehensum. P Et refertur vtrumque ad pronomen te. Et est sensus. praecipio tibi te imaculatumlirrepraehensum seruare mandatum. Non sufficit mihi vt serues mandatum, sed exigo vt imaculatus intus et irrepraẽhensibilis foris serues mandatum. Serues autem non ad diem aut annum, sed VSQVE in aduentum. pro. apparitionendñi nostri iesu christi. ] hoc est vsque  in finemjrationem habendo appartionis domini nostri iesu christi ad indicandum omnes et scrutandum singula hominum opera. QVEM. pro. quam. P Refert enim apparirionem. SVIS temporibus.per soli deo notis. OSTENDET beatus. P Ad apprçhensionem vitae aeternae inuitans aptedescribit deum beatum. Rursus ap paritionem christi ad iudicium ad beandum electos describens ab autore deojaptissime noiat deum bearum, vt ppriam camm ois beatitudinis tangat. ET solus potens.P Vere solus pont:cuius potestas ad oia se ex tendit, et a quo omnium potestates sic pendent vtsi velit nihil possint. Et vtrumque horum (Lesse beatum et esse solum potentem) naturale est deo. REX regum. pro. regnantisii et dominus dominantium.) Vtrumque horum ad gubemationem spectat vniuersi. QVI solus habet imortalitatem. P Omnis mutatiojquaedam morris soecies est et perinde dictum est quod solus habet imortalitaremjac si dictum fuisset qui solus viuit absque oi mutabilitatenam omnis creatura quantuncunque imortalis imo etiam beatalalicui subiacet mutationi.nam et de facto motus mentis aliquos habet successiue: praeter hoc quod mutabilis est per porentiam diuinam. ET LV CEM inhabitat, pro. inhabitans inaccessibilem, p Superfluit et. De luce spiritu ali est sermo. Et diuina substantia appellatur intelligibilis lux inaccessibilis tantae excellentiae vt nullus possit ad eam accedere. Quemadmodum lur solis est inaccessibilis relatine ad corporea, ita lux diuina est inaccessibilis relatiue ad totum vniuersum. Metaphorico autem sermone dicit quod inhabitat. nulla siquidem est ibi proprie inhabiratio:sed ad similitudi nem habitantis in luce inaccessibiliiexcellentiam diuinae lucis hoibus describit. QVEM nullus hoïum vidit: sed nec videri pro. videreipont P Ex hac secunda particula apparet quod de hoie fm naturae vires loquit. nam sic absque omni controuerlia verificatur quod nullus hoium deum vidir et quod nec videre eum pont. Et hoc intellige de videre deum fm seipsum: ad differentiam eius quod est videre deum in aliossine intelligibili fiue imaginabili siue sensibili. nam sic pphetae viderunt deum: et philosophi et angeli naturaliter vident deuĩ, CVI honor et imperium sempiternum:amen. P Par fuit vt descripto deo fm tot praeclaras conditiones, in gratiarumactiones prorumperet: cupiendo exhiberi illi honorem et aduenire imperium sempitemnum quemadmodum in coelo ita et in terra. Possunt tamen et haec assertiue intelligi:vt etiam describat deum c honolall et imperare sempiterno tempore.  \pend\pstart ¶’DIVITIBVS huius saeculi praecipe non sublime sapere. pro. senrire. P Necessaria monstrat haec esse dit cendo praecipe. Et vere primum praeceptum negariuum quo egent diuiresjest non alte sentire. Diuitias eni magnipendunt: ac perhoc efferunt animum magnipendentium illas et possidentium. Elata ergo mens primo inhibetur deinde. NEQVE sperare in incerto. pro. in incertitudinejdiuiriars̃ a Et hoc quoque familiare est diuiribus vt confidant in auro. Et vere dicit in incertitudine diuitiarti: quoniam incerta est earum stabilitas incertus est earum futurus vsus et incerrior fructus. SED IN deo viuo, qui prestat nobis omnia abunde ad fruendum.per hoc est ad delectabiliter suauiterque vtendum. Et quod dicit abundesntellige quantum est ex parte dei nam secus quandoque accidit ex parte nostri: vel pprer demerita nostraivel propter aliquod spirituale bonum. BENE agere. Refertur ad praecipe diutibus huius saeculi. Nec dicit bona sed bene agere, vt modum agendi ex parte agentis explicet bonum DIVITES fieri in operibus bonis mul\pend
\marginpar{[ p.147 ]}
\section*{.II. AD TIMO. }
\section*{CAPVT .I. }\pstart tiplicando bona opera. FACILE tribuere.p non solum tribuere sedfacile, COMMVNICARE. P Ra tionem facile tribuendi explicar adiungendo communicationem. Diuires siquidem vrpore abundantes bo nis fortunae considerare debent q ideo fm ordinem dininae prouidentiae abundant vt communicent. non enim eis est data abundantia vt tot propria rerineant, sed vt propria quibus abundant communia efficiant. et propterea faciles debent esse ad tribuendum. Er ne hoc reputarent esse sibi damnumssubiungit. THESAVRIZARE sibi fundamentum bonum in futurum. per Ecce fructus eommunicandi. ipsum siquidem con municare est recondere sibiipsis bonum fundamentum in futurum. Et ad ad sit hoc fundamentum subiungir. VT appraehendant veram. pro.aeternam, vitam. p Electionis recondit fundamenrum super quod consurget ascensus ad appraehendendum vitam aeternam, hoc est. n. pprium meritum diuitum quatenus sunt diuites. LO TIMOTHEE.P Ad excitandum timotheumjadiungit articulum. DEPOSITVM custodi.] Con missum gregem sub depositi rarione custodiendum mandat:vt eriam de leui culpa teneatur. DEVITANS prophanas vocum nouitates, pro. inanitates. P Inter alias vigilantias necessarias ad custodiam depositii oportet te euitare prophanas vocum inanitates. Inanitares vocum aliae sunt prophanaejhoc est sacris aduersaetaliae sunt tanq̃ indifferentesivt nugae. et propterea paulus addidit prophanas. Oportet siquidem episcopum custodire depositum sibi gregem eriam a vanis vocibus si prophanae suntisi aduersus sacras literas suntsi aduersus sacta mysteria sunt. ET oppositiones falsi nominis scientiae.per hoc est falso nominatae scientiae. Forte scientiam cabalisticam appellat falso nominatam scientiam. verum quancunque intelligas, perpende soliditatem veritatis cuangelicae nam non contrariantur ei oppositiones cuiuscumque scientiaeidummodo non falso nominetur scientia:sed solum opposiriones falso nominatae scientiaejadnersantur veritati christianae. Et in promptu ratio est:quia verum vero non est contrarium. omnis autem vere nominata scientia est verorum. QVAM quidam promittentes, pro. profitentes circa fidem exciderunt. pro. a fide aberrauerunt.P Esperientiam affert quo ducit illa falso nominata scientia. Et hinc habemus quod de quadam certa scientia loquiturtalioquin non diceret quam quidam profitentes a fide aberrauerunt. Et si collegeris hec omnia (videlicet quod falso nominat scientiamque complexa est a fidelibus cum periculo profitentium illam ita quod quidam illam profitentes aberrarunt a fide: et quod paulus non dicit illam omnino vitandanssed dicit illius virandas esse duas partes.sprophanas vacuitates vocum et opposiriones ad custodiendum depositum gregem) suspicaberis forte mecum cabalisticam pingi scientiam. Meae siquidem suspitioni alia non occurrit scientia his picta colonbus. GRATIA tecum amen. P Haec est salutatio manu pauli.P Et apud graecos inuenit subsenptio talis. Ad timotheum prima scripta laodiciç̨. Caietae die vltimo martii.1529.  \pend
\phantomsection
\addcontentsline{toc}{subsection}{\textit{THOMAE DE VIOCAIETAT NLCARDINALIS SANCTI XISTI:IN POSTERIOREM EPISTOLAM PAVLI AD TIMOTHEVM COMMENTARII. }}
\subsection*{\textit{THOMAE DE VIOCAIETAT NLCARDINALIS SANCTI XISTI:IN POSTERIOREM EPISTOLAM PAVLI AD TIMOTHEVM COMMENTARII. }}\pstart \huge\textbf{P}\normalsize AVLVS apostolus iesu christi.per non aliog apostolor. PER volun tatem deiper non meritis propriis. SECVNDVM promissionem vitae quae est in chnsto iesu.] apostolus per voluntarem dei non fm quod disponit et gu bernat naturalem rerum ordinem, sed fm quod pmittit vitam quae est in christo iesu nune per fidemjpostremo autem per imortalem gloriam enam corpo ris. Materia siquidem apostolatus significatur: explicando quod est missus per voluntatem dei quatenus promittit vitam quae est in christo iesu. ad hanc siquidem promissionem promulgandamjad illam inuitandum etc. legatus est. TIMOTHEO charissimo, pro. dilectoifilio.) per euangelium. GRATIA et misericordia et pax.P Superfluunt ambae coniunctiones. Imprecatur non solum extrema donacgratiam et pacem) sed etiam misericordiamjdonum medium summe necessanum praelatis. A DEO patre nostro.P Superduit nostro. ET CHRISTO Iesu domino nostro.P Exposita saepesunt haec.  \pend\pstart \huge\textbf{G}\normalsize RATIAS ago. pro, gratiam habeojdeo meo.P Superfluit meo. Perfecta salutatione inchoat narrationem a commendatione timorhei. Er commendacionem inchoat explicando se recognoscere vt diuinum beneficium quod assidueest memor ipsius timothei, memor aurem in orarionibus et in desiderio videm di ob virtutes eius et mattis et auiae. CVI seruio a progenitoribus meis in conscientia pura ) Superfluit Caiet. super epi. apo.  \pend
\endnumbering
\end{pages}
\end{document}
        